\providecommand{\phio}{\phi(0)}
\providecommand{\cR}{\mathcal R}
\providecommand{\Prb}{\mathbb P}
\providecommand{\nm}{\mathrm}
\section{Acquisition and comparison near a two-point law}
\label{j:local:section}
The defect identity has two uses. Contact and a complete Fourier bound
first confine one extremizing marginal to a neighborhood of a two-point
law. We then compare the actual sum with that law's own binomial count.
The second step needs no extremality beyond the simultaneous hypotheses
already obtained in the first. All bounds below include zero excess
$e=C_n-\CE=0$; positive conditional noise and zero noise are treated
separately at the end.

Write $\sigma=\sqrt{10}$, $m=\E X_+=\E X_-$,
$\rho=\E|X|^3$, $\gamma=\E X^3$, $\beta=\E[X|X|]$,
$\mu_j=\E X^j$ and $M_j=\E|X|^j$ (also $\eta_5=M_5$).
In the central gap we use $(-A,B)$ for the endpoints, equivalently
$A=A_-$ and $B=A_+$ in Section~\ref{j:core:section}.
The quantities $J,J_1,J_2,D,L,K_1,K_q$ and $\mathcal R$ are those of
Theorem~\ref{j:core:conservation}; $U=J_1+(\mu_4-3)/6$ below is a
coefficient, distinct from the conditional noise introduced later.
The letter $B$ denotes the positive endpoint until Bernoulli coordinates
are introduced; there its meaning is explicitly reset to a label.


\subsection{Support from the contact inequality}\label{j:local:support:theorem}
\paragraph{The scale correction.}
Write $\phi_0=\phi(0)$.
For $n>1$ put $s_n=\sqrt{n/(n-1)}$ and
$e_n=n(s_n-1)-1/2$.
\begin{lemma}\label{j:local:support:scale}
For every integer $n\ge36$, one has $0<e_n<.384/n$.
\end{lemma}
\begin{proof}
Put $x=1/n$ and $c=.384=48/125$. Direct expansion gives
\[
 (1-x)(1+x/2+cx^2)^2-1
 =x^2\{(2c-3/4)-(c+1/4)x+(c^2-c)x^2-c^2x^3\}.
\]
The bracket is decreasing for $x\ge0$, since $0<c<1$.
At $x=1/36$ its exact rational value is $823/4050000>0$.
Squaring positive quantities gives
$s_n<1+x/2+cx^2$ for $0<x\le1/36$. Similarly
$(1-x)(1+x/2)^2-1=-3x^2/4-x^3/4<0$, which gives the lower bound.
\end{proof}


\paragraph{A signed Gaussian envelope.}
Define, for every real $H,t$ and $n>1$,
\[
 \mathcal H_n(H,t)=n^{3/2}\{\Phi((H-t)/\sqrt{n-1})-
             \Phi(H/\sqrt n)\}
       +nt\phi(H/\sqrt n)+\tfrac H2(t^2-1)\phi(H/\sqrt n).
\]
Write $v=n-1$, $s=\sqrt{n/v}$, $u=|t|$, and
\[
 A=\frac{u^3-3u}{6},\quad B=\frac{u^4-6u^2+3}{24\sqrt v},
 \quad D=\frac{u^5-15u}{120v}.
\]
\begin{lemma}[Fifth-order identity with a uniform finite remainder]
If $u>\sqrt3$, then, with $z=H/\sqrt v$,
\[
 \frac{\mathcal H_n(H,\pm u)}{s^3}
 =\pm A(1-z^2)\phi(z)+B\phi^{(3)}(z)
                  \mp D\phi^{(4)}(z)+R,
\]
where the sign of the remainder need not be fixed and
\[
 |R|\le E^{(5)}_v(u):=
 \frac{51u^6}{14400v^{3/2}}
 +\frac{51(3u^2-2)}{960v^{3/2}}+\frac{u}{8v^2}.
\]
This statement is uniform over every real $H$.
\end{lemma}
\begin{proof}
Let $F_w(x)=\Phi(x/\sqrt w)$, with primes denoting $x$ derivatives.
The heat identity $\partial_w F_w=F_w''/2$ and ordinary Taylor's theorem
give, with unlabelled derivatives at $(v,H)$,
\begin{align*}
 F_{v+1}&=F_v+F_v''/2+F_v^{(4)}/8+R_0,
 &|R_0|&\le\sup_{w\in[v,v+1]}\|F_w^{(6)}\|_\infty/48,\\
 F_{v+1}'&=F_v'+F_v'''/2+F_v^{(5)}/8+R_1,
 &|R_1|&\le\sup_{w\in[v,v+1]}\|F_w^{(7)}\|_\infty/48,\\
 F_{v+1}''&=F_v''+F_v^{(4)}/2+R_2,
 &|R_2|&\le\sup_{w\in[v,v+1]}\|F_w^{(6)}\|_\infty/8.
\end{align*}
Expand $F_v(H-t)$ through degree five, with remainder at most
$u^6\|F_v^{(6)}\|_\infty/720$. In
\[
 \frac{\mathcal H_n}{n^{3/2}}=
 F_v(H-t)-F_{v+1}(H)+tF_{v+1}'(H)
             -\frac{t^2-1}{2}F_{v+1}''(H),
\]
the first three derivative coefficients cancel. The next three are
\[
 \left(\frac t2-\frac{t^3}6\right)F_v'''
 +\frac{t^4-6t^2+3}{24}F_v^{(4)}
 +\left(\frac t8-\frac{t^5}{120}\right)F_v^{(5)}.
\]
The remaining absolute value is bounded by
\[
 \left\{\frac{u^6}{720}+\frac1{48}+\frac{u^2-1}{16}\right\}
       \sup_{w\in[v,v+1]}\|F_w^{(6)}\|_\infty
 +\frac u{48}\sup_{w\in[v,v+1]}\|F_w^{(7)}\|_\infty.
\]
Fourier inversion of the Gaussian, differentiated under its absolutely
integrable Fourier transform, gives
\[
 \|\phi^{(k)}\|_\infty
 \le\frac1\pi\int_0^\infty y^k e^{-y^2/2}\,dy.
\]
Integration by parts yields
$\|\phi^{(5)}\|_\infty\le8/\pi<2.55$ and
$\|\phi^{(6)}\|_\infty\le15\phi_0<6$.
Since $F_w^{(j)}=w^{-j/2}\phi^{(j-1)}$, the displayed remainder
becomes $s^3E^{(5)}_v(u)$ after multiplication by $n^{3/2}$.
The three retained derivative coefficients give the identity.
\end{proof}

\begin{lemma}[A lower tangent bound for the fourth Gaussian derivative]
For every real $z$,
\[
 \phi^{(4)}(z)\ge\phi_0(3-15z^2/2).
\]
\end{lemma}
\begin{proof}
The same Gaussian Fourier integral gives
$\phi^{(4)}(z)=\pi^{-1}\int_0^\infty
y^4e^{-y^2/2}\cos(yz)\,dy$.
Use the globally valid inequality $\cos w\ge1-w^2/2$ inside
this integral. Its fourth and sixth Gaussian moments are $3\phi_0$
and $15\phi_0$, respectively.
\end{proof}

\begin{lemma}[Uniform signed fifth-order envelope]\label{j:local:support:envelope}
Suppose $u>\sqrt3$ and $A,B,D\ge0$. If
$W=A-15D/2>0$, put
\[
 T=A-3D+\frac{9B^2}{4W}.
\]
If $T\ge A/2+2B+3D$, then, for every real $H$ and either sign,
\[
 \mathcal H_n(H,\pm u)\le s^3\{\phi_0 T+E^{(5)}_v(u)\}.
\]
\end{lemma}
\begin{proof}
For positive $t=u$ and $|z|\le1$, divide the retained terms by
$\phi_0$ and use the preceding tangent bound to obtain
\[
 A-3D-(A-15D/2)z^2+3B|z|
 \le A-3D+9B^2/(4W).
\]
Here $(1-z^2)\phi(z)\le(1-z^2)\phi_0$ and
$|\phi^{(3)}(z)|\le3|z|\phi_0$ on this interval.
For $|z|\ge1$ the first term is nonpositive, and the Fourier bounds
$\|\phi^{(3)}\|_\infty<2\phi_0$ and
$\|\phi^{(4)}\|_\infty\le3\phi_0$ give $2B+3D$.
For negative $t=-u$, the first term has global upper bound
$2e^{-3/2}A\phi_0<A\phi_0/2$, and the other terms give
$(2B+3D)\phi_0$. The hypothesis on $T$ covers both exterior cases.
The inequality $2e^{-3/2}<1/2$ follows from four positive terms of
the exponential series at $3/2$. Finally add the uniform remainder.
\end{proof}

\begin{lemma}[A global Gaussian remainder bound]
\label{j:local:ext:global-gaussian-ray}
For every integer $n\ge2$ and all real $H,t$,
\[
 |\mathcal H_n(H,t)|
 <\frac{|t|^3}{15}+\frac{|t|}{5}+\frac14
 <\frac{|t|^3}{15}+\frac{|t|}{5}+\frac12.
\]
\end{lemma}
\begin{proof}
Write $F_w(x)=\Phi(x/\sqrt w)$, $v=n-1$ and $y=H-t$.
The identity defining $\mathcal H_n$ is
\[
 n^{-3/2}\mathcal H_n(H,t)
 =F_v(y)-F_n(H)+tF_n'(H)-\frac{t^2-1}{2}F_n''(H).
\]
Taylor expansion of $F_n(y)$ at $H$ through degree two has remainder
at most $\phi(0)|t|^3/(6n^{3/2})$, since
$\|\phi''\|_\infty=\phi(0)$.
After this expansion the remaining correction is
\[
 \frac12\{F_n''(H)-F_n''(y)\}
 +\frac12\int_v^n\{F_n''(y)-F_w''(y)\}\,dw.
\]
Here we used $\partial_wF_w=F_w''/2$. The first term is at most
$\phi(0)|t|/(2n^{3/2})$ in absolute value. Applying the heat identity
once more bounds the second term by
\[
 \frac14\int_v^n\int_w^n
       |F_r^{(4)}(y)|\,dr\,dw
 \le\frac{\|\phi'''\|_\infty}{8v^2}
 \le\frac1{4\pi v^2}.
\]
The last inequality follows by differentiating the absolutely
convergent Gaussian Fourier integral, which gives
$\|\phi'''\|_\infty\le2/\pi$. Multiply by $n^{3/2}$ and use
$\phi(0)<2/5$. The function $x^{3/2}/(x-1)^2$ decreases for $x>1$,
because its logarithmic derivative is
$3/(2x)-2/(x-1)=-(x+3)/(2x(x-1))<0$. Hence, for $n\ge2$,
\[
 \frac{n^{3/2}}{4\pi(n-1)^2}
 \le\frac{\sqrt2}{2\pi}<\frac14,
\]
using $\sqrt2<3/2$ and $\pi>3$. This proves the stronger bound and
therefore the exact ray majorant used by the original checkers.
\end{proof}

\paragraph{Finite bands and the exact support polynomial.}\label{j:local:support:polynomial}
For integers $36\le N\le L$, use the simultaneous endpoint bounds
\[
 B_+=\frac{u^4-6u^2+3}{24\sqrt{N-1}},\quad
 D_+=\frac{u^5-15u}{120(N-1)},\quad
 D_-=\frac{u^5-15u}{120(L-1)}.
\]
For an infinite index range, set $D_-=0$ instead. Put
\[
 W_*=A-15D_+/2,\qquad
 T_*=A-3D_-+\frac{9B_+^2}{4W_*}.
\]
The proof of Lemma~\ref{j:local:support:envelope} applies uniformly on the whole band
if $W_*>0$ and $T_*\ge A/2+2B_++3D_+$. In particular the favorable
term is frozen at the upper index endpoint, and the adverse terms at
the lower one. No monotonicity of a signed coefficient is assumed.

Let $(P,H)$ be an attained positive closed maximizing pair, standardized
to mean zero and variance one. Write $C=C_n\ge C_E$,
$\rho=\mathbb E|X|^3$, $\beta=\mathbb E[X|X|]$, and
$d=n(C_{n-1}-C_n)_+$. The finite contact/predecessor bound is
\[
 \rho I(t)\le\mathcal H_n(H,t)+
 C\{-|t|^3+3\rho t^2/2+3\beta t+\rho e_n\}+\rho s_nd,
\]
with $I=0$ on the entire closed support. Suppose also, for a finite
band, that $C_n-C_E\ge K/n$, where $K>0$. The finite moment bounds of~\cite[p.~125]{Shevtsova2013}, namely
$C\rho\le .3031(\rho+.646)$ and
$C\rho\le .3322(\rho+.429)$, give
\[
 c=.4095+K/L<C_n,\qquad
 \rho<r=\min\left\{\frac{.3031(.646)}{c-.3031},
                        \frac{.3322(.429)}{c-.3322}\right\}.
\]
For the unrestricted nonnegative case put $K=0$, $c=.4095$.
The sign-moment inequality
$\rho\ge(1+\beta^2)/\sqrt{1-\beta^2}$ supplies $|\beta|<b$ whenever
$(1+b^2)^2>r^2(1-b^2)$ and $0<b<1$. All these statements concern the
same current law. The lower endpoint cannot replace $L$ in $K/L$.

When a preceding resource enclosure supplies sharper actual bounds
$\rho\le r$ and $|\beta|\le b$, these may instead be used in the
same polynomials. Their validity is the only extra premise.

Use rational $\ell^2\le N-1$ and $S^2>N/(N-1)$. For $u=|t|$ define
\[
 Q(u)=-u^3+3ru^2/2+3bu+.384r/N.
\]
First check $Q<0$, before replacing $C$ by its lower bound $c$.
Replace $1/\sqrt{N-1}$ by $1/\ell$ in the adverse remainder and in
the term $2B_+$ of the domination gate. Write $E_*$ for that rational
remainder. Strict negative contact for every $n$ in the band follows
from the polynomial inequality
\[
 \frac25S^3\left\{(A-3D_-)W_*+
              \frac{(u^4-6u^2+3)^2}{256(N-1)}\right\}
 +S^3E_*W_*+\{cQ+rSq\}W_*<0.
\]
The second polynomial gate is $W_*>0$. The third is
\[
 \{A/2-3D_--3D_+-2B_+\}W_*+
         \frac{(u^4-6u^2+3)^2}{256(N-1)}>0.
\]
This is the cleared-denominator version of the domination hypothesis.
The bound on $T_*$ is positive, so the replacement $\phi_0<.4$ and
$s_n<S$ both have their correct direction.

The exact kernel tests these polynomials on all adjacent closed intervals
of length $.05$ from $R$ to $4.5$. Positive Bernstein coefficients prove
their signs on every point of each cell, including endpoints. On the
rays $u\ge4.5$ use Lemma~\ref{j:local:ext:global-gaussian-ray}, which gives the same bound
$|\mathcal H_n(H,t)|\le u^3/15+u/5+1/2$.
The resulting cubic and its first three derivatives are strictly negative
at $4.5$, as is the coefficient $Q$ before replacing $C$ by $c$.
This excludes the whole rays. Thus, if $d\le q$, actual support is in
$(-R,R)$, including exclusion of arbitrarily small exterior atoms and
exterior accumulation points.


\paragraph{The support rows used in this proof.}
For the all-index local argument, take $K=0$, $R=3.84$, $q=.19$ on
\[
[45,49],\ [50,59],\ [60,89],\ [90,134],\ [135,199],\ [200,\infty).
\]
Only indices at least 47 are used. The last row uses $D_-=0$;
there is no upper index restriction. The exact checker reconstructs
all closed $u$ cells and the two exterior rays. Missing an index band
or the infinite row is an acceptance failure.

\paragraph{One bounded-threshold series.}\label{j:local:support:series}
Put $c_k=(-1)^k/[2^kk!(2k+1)]$, and let $(x)_j$ denote the rising
factorial. Define the following rational polynomials:
\begin{align*}
 P_j(H,t)={}&\sum_{k=0}^{j+1}
 \frac{c_k(k+1/2)_{j+1-k}}{(j+1-k)!}(H-t)^{2k+1}
       -c_{j+1}H^{2j+3}\\
 &+\frac{(-1)^{j+1}tH^{2j+2}}{2^{j+1}(j+1)!}
       +\frac{(-1)^jH^{2j+1}(t^2-1)}{2^{j+1}j!}.
\end{align*}
In particular,
\[
 P_0=(t^3-3t)/6,\qquad
 P_1=\frac{3t-t^3}{4}H^2+\frac{t^4-6t^2+3}{8}H
                  +\frac{-t^5+10t^3-15t}{40}.
\]
Fix $h>0$, $E(x)=(193/71)^{\lceil x\rceil}$, and
\[
 B_r=\frac{7(h+r)}6E((h+r)^2/6)
       +\left(h+r+\frac{h(r^2+1)}8\right)E(h^2/8).
\]
\begin{lemma}\label{j:local:support:finite-series}
For every $n\ge4$, $|H|\le h$ and $|t|\le r$,
\[
 \left|\mathcal H_n/\phi(0)-\sum_{j=0}^5P_j/n^j\right|
       \le4^7B_r/n^6.
\]
\end{lemma}
\begin{proof}
Put $z=1/n$. Integration of the Gaussian exponential series gives
$\mathcal H_n/\phi(0)=S(z)/z$, where
\begin{align*}
 S(z)={}&\sum_{k\ge0}c_k\{(H-t)^{2k+1}(1-z)^{-k-1/2}
                                  -H^{2k+1}\}z^k\\
 &+te^{-H^2z/2}+\tfrac12zH(t^2-1)e^{-H^2z/2}.
\end{align*}
Its constant coefficient vanishes, and expansion of the binomial factor
gives $[z^{j+1}]S=P_j$. At $z_0=1/4$, the sum obtained by replacing
all coefficients by their absolute-value majorants is at most $B_r$:
use $(1-z_0)^{-1/2}<7/6$, $1/(2k+1)\le1$, and $e<193/71$.
The latter bound follows, for example, by summing the exponential series
through degree six and bounding the rest by a geometric series with
ratio $1/8$. All products and rearrangements are therefore absolutely
convergent. The discarded terms have degrees at least seven in $S$,
so their sum at $0\le z\le z_0$ is at most $(z/z_0)^7B_r$.
Divide by $z$ to get the result. This proves a finite uniform bound,
including every threshold in the stated interval and every $n\ge4$.
\end{proof}

\paragraph{A common contact test for radius and strip refinement.}
\label{j:local:support:panels}
Let $N\ge36$, $E,q\ge0$ and $C_0>0$. Suppose $n\ge N$,
$0\le C-C_E\le E/n$, $\delta_n:=n(C_{n-1}-C_n)_+\le q$, $C\le C_0$ and
$|X|\le R_0$. A previously acquired row gives
\[
 \rho\le r,\quad 0\le b_l\le\beta\le b_u<1,\quad |H|\le h.
\]
Define $Q(t)=-|t|^3+3\rho t^2/2+3\beta t$ and
$L_P(t)=C_EQ(t)+\phi(0)(t^3-3t)/6$.
For $t>0$, upper polynomials for $L_P(t),Q(t)$ are
\[
 (1/15-.4095)t^3+1.5(.40974)rt^2+(3(.40974)b_u-.3988/2)t,
 \quad -t^3+1.5rt^2+3b_ut.
\]
For $t<0$, write $u=|t|$ and use
\[
 (-.3988/6-.4095)u^3+1.5(.40974)ru^2+(.2-3(.4095)b_l)u,
 \quad -u^3+1.5ru^2-3b_lu.
\]
These coefficient choices are valid because $b_l\ge0$ is checked first.
On any closed cell of one sign, let $L_J$ be an upper bound for the
first polynomial and let $Q_J\ge0$ bound the positive part of the second.
For $j=1,2,3$, take $m_j$ to be the positive part of the largest tensor
Bernstein coefficient of $P_j$ on $[-h,h]\times J$. For $j=4,5$, the
absolute monomial bound
$\sum_{a,b}|[H^at^b]P_j|h^aR_J^b$ suffices, where
$R_J=\max_{t\in J}|t|$. The support contact is then bounded above by
\begin{equation}\label{j:local:support:paneltest}
 L_J+\frac{E Q_J}{N}
 +.4\left(\sum_{j=1}^5\frac{m_j}{N^j}
                +\frac{4^7B_{R_J}}{N^6}\right)
 +r\left(qS+\frac{.384C_0}{N}\right),
 \qquad S>\sqrt{N/(N-1)}.
\end{equation}
Indeed substitute the finite series in contact, split $C=C_E+(C-C_E)$,
and use the scale estimate $e_n<.384/n$. Every error retained in
\eqref{j:local:support:paneltest} is nonnegative and decreases with $n$.
A strictly negative rational result therefore excludes the entire cell
for every integer $n\ge N$, including atoms and accumulation points.
The uniform-tail specialization is $(N,E,q,C_0)=(47,.7,.19,.425)$;
its score cap follows from $C<.40974+.7/47<.425$.
The finite bands use their own acquired $E,q,C_0$ at the corresponding
floor, with exactly the same formula.

For a coordinate strip, conditional centering gives
\[
 X=\ell(Y-p),\quad
 \ell=\sqrt{(1-\theta)/(p(1-p))},\quad
 \beta=(1-\theta)(1-2p)+\kappa,\quad |\kappa|\le\theta.
\]
If $p\in[p_l,p_u]\subset(0,1/2)$ and $\theta\le\Theta$, put
$\ell_l\le\sqrt{(1-\Theta)/(p_u(1-p_u))}$.
The forbidden strip $Y\notin(-.55,1.55)$ maps into
\[
 X\le-\ell_l(p_l+.55),\qquad X\ge\ell_l(1.55-p_u).
\]
Apply \eqref{j:local:support:paneltest} between these thresholds and $\pm R_0$.
An interval beyond the known radius is empty; an equal endpoint is a
singleton that is still checked. If only an upper $\beta$ bound is given,
use $b_l=(1-\Theta)(1-2p_u)-\Theta$, after checking its nonnegativity.

The selected uniform consequences used in the resource schedule are:
\[
\begin{array}{l|c|c|r}
\text{input rows}&\text{prior radius}&\text{conclusion}&\text{closed cells}\\\hline
10\text{ central preparation rows}&3.84&|X|<2.5&1080\\
45\text{ central rows}&2.5&-.55<Y<1.55&1431
\end{array}
\]
The respective largest contact bounds are below $-.084$ and $-.097$.
These implications acquire no row by themselves: each is used only after
its input row is proved by the preceding resource step.
For the raw left branch, the all-threshold support polynomial with
$(N,L,K,q)=(47,\infty,0,.19)$ and actual moment caps
$(r,b,R)=(1.390,.430,3)$ gives radius $3$.
The same heat proof supplies the initial radius and the finite-base
support rows, with their stated parameters. The bounded-threshold series
supplies the central contraction and coordinate strip above, and the
separately specified earlier finite strips.


\subsection{Acquiring simultaneous resources for one marginal}\label{j:local:sign}
For a ratio maximizer $C=C_n\ge\CE$, put $e=C-\CE$ and $w=H^2/(2n)$.
The core identity is
\[
 ne\rho+n\CE\rho(1-e^{-w})+n\phi(0)e^{-w}D/6+n^{3/2}L=\mathcal R.
\]
The first complete raw enclosure must give $\mathcal R<r<.4N$, for the
actual $n\ge N$. Since $C\rho\ge\CE>.4$, this licenses
Corollary~\ref{j:core:license}, including
\[
 ne\rho+n\phi(0)D/6+n^{3/2}L\le\mathcal R,\quad
 e^{-w}\ge1-5r/(2n),\quad nD\le15000r/997.
\]
The proof uses $e\ge0$ and never divides by $e$.
If a previous pass gives $nD\le Nd$, subtracting the favorable phase
term gives the sharper shared budget
\begin{equation}\label{j:local:budget}
 ne\rho+\kappa nD+n^{3/2}L\le\mathcal B(H),\qquad\kappa=997/15000,
\end{equation}
by increasing each raw quadratic curvature by
\begin{equation}\label{j:local:re:extra}
 q_{\rm extra}(\rho)=\tfrac12(1-5r/(2N))(.4095\rho-d/15)>0.
\end{equation}
Indeed $1-e^{-w}\ge we^{-w}$, $D\le d$, $\phi(0)<.4$ and
$\CE>.4095$. The positivity gate is checked before this use.
The field called \texttt{Qbudget} already contains this addition.
For a cell with $D\ge d_l\ge0$ and $\rho\ge\rho_l>0$, the same identity
also gives the score update
\[
 C\le .40974+\{r/N-.3988(1-5r/(2N))d_l/6\}/\rho_l.
\]
A negative numerator excludes the cell by $e\ge0$. Otherwise the numerator
increases with $1/n$, proving this bound at every $n\ge N$.


\paragraph{One positive-measure rule.}
Suppose zero lies in the gap $(-A,B)$, $A,B>0$. Let
\[\begin{gathered}
 p=\Pp(X>0),\quad q=1-p,\quad m=\E X_+=\E X_-,\quad V=\E[X\mid\mathbf1_{X>0}],\\
 W=X-V,\quad\theta_\pm=\E[W^2;\pm X>0],\quad\theta=\theta_++\theta_-.
\end{gathered}\]
Then $m^2=(1-\theta)pq$, and the conditional means have magnitudes
$\bar A=m/q$, $\bar B=m/p$. Set
$\Delta_A=\bar A-A\ge0$, $\Delta_B=\bar B-B\ge0$.
We also use the unnormalised endpoint displacements
\[
 M_+=\E[(X-B)\mathbf1_{X>0}]=p\Delta_B,\qquad
 M_-=\E[(-X-A)\mathbf1_{X<0}]=q\Delta_A.
\]

\begin{lemma}[Positive-measure pricing]\label{j:local:re:pricing}
For $x=|X|$, $a=|V|$, and a function $h$ constant on each sign class,
\begin{equation}\label{j:local:re:conditional}
 \E[h(|X|^j-|V|^j)]
 =\E\left[hW^2\sum_{k=0}^{j-2}(k+1)a^kx^{j-2-k}\right],\quad j\ge2.
\end{equation}
Also, for $d\mu_0=(|X|-1)^2(|X|+1/2)\,dP$,
\begin{align*}
 \mu_0(\R)&=\rho-1,\\
 M_4-1&=\int\frac{(x+1)^2}{x+1/2}\,d\mu_0,\\
 M_5-1&=\int\frac{x^3+2x^2+3x+3/2}{x+1/2}\,d\mu_0,\\
 K_1+3\rho\beta&=\int
  \frac{\operatorname{sgn}(X)x(x+2-3\rho/2)}{x+1/2}\,d\mu_0.
\end{align*}
For any nonnegative resources $z_i$ satisfying $\sum a_i z_i\le r$,
$a_i>0$, a signed linear cost obeys
\begin{equation}\label{j:local:re:simplex}
 \sum b_i z_i\le r\max\{0,b_1/a_1,\ldots,b_k/a_k\}.
\end{equation}
\end{lemma}
\begin{proof}
Divide $x^j-a^j-ja^{j-1}(x-a)$ by $(x-a)^2$; the linear term has
zero conditional mean. The radial identities follow by multiplying
out and using $\E X=0$, $\E X^2=1$. Finally each $b_i$ is at most
$a_i$ times the displayed maximum, and all $z_i$ are nonnegative.
\end{proof}

This one rule accounts for the moment, signed Taylor, endpoint movement,
and outer-gap estimates. In particular, no two costs
may each spend the full residual and then be added.



The gap identity and simultaneous defect budget are
\eqref{j:core:gap-defect} and \eqref{j:core:simultaneous}, with
$(A_-,A_+)=(A,B)$ and $(\Delta_-,\Delta_+)=(\Delta_A,\Delta_B)$.
The outer-gap estimate \eqref{j:core:outer-gaps} bounds the same
$T_\pm,R_\pm,J_o$. One resource pays for all their charges together.

\paragraph{A state, its update, and its exported resource.}
A state for a closed probability interval $I=[l,u]$ records
\[
 |X|\le R,\quad C\in[c_0,c_1],\quad\mathcal R\le r,
 \quad nD\le Nd,\quad\theta\le\Theta<1,
 \quad A\ge A_l>0,\quad B\ge B_l>0,
\]
together with simultaneous signed-moment and coefficient intervals.
The distinction between $nD\le Nd$ and $D\le d$ is essential.
Let $s_+=\max_Ip(1-p)$, $s_-=\min_Ip(1-p)$ and
\begin{gather*}
 b_I=\frac{(6+2\sigma)(1-\Theta)^{3/2}}{\sqrt{s_+}}
                 \operatorname{dist}(p_*,I)^2,
 \qquad r_I=d-b_I,\\
 K_+=3(\sigma+1)B_l,\quad K_-=3(\sigma+3)A_l,
 \quad K=\min(K_+,K_-),\quad m_l=\sqrt{(1-\Theta)s_-}.
\end{gather*}
If $r_I<0$, exclude the entire closed cell. Otherwise \eqref{j:core:simultaneous}
gives, simultaneously,
\begin{equation}\label{j:local:re:residual}
 6m(A_l+B_l)(\Delta_A+\Delta_B)+K_+\theta_++K_-\theta_-\le r_I.
\end{equation}
Thus the actual law satisfies the intersected bounds
\begin{gather}
 \theta\le\min\{\Theta,r_I/K\},\notag\\
 A\ge\max\left\{A_l,\sqrt{(1-\Theta)l/(1-l)}-
                  \frac{r_I}{6m_l(A_l+B_l)}\right\},\notag\\
 B\ge\max\left\{B_l,\sqrt{(1-\Theta)(1-u)/u}-
                  \frac{r_I}{6m_l(A_l+B_l)}\right\}.\label{j:local:re:geometry}
\end{gather}
Each right side is evaluated from established bounds; sequential
improvements may then be used by the next expression. A nonpositive
endpoint estimate fails this update, rather than excluding a law.
Initialization uses the same formulas with the valid lower sum
$A+B\ge\sqrt{2-d}$.

In Bernoulli coordinates set $Y=B_0+U$, $\E[U\mid B_0]=0$,
$B_0\sim\mathrm{Bernoulli}(p)$, $X=\ell(Y-p)$ and
$d_Y=(\Delta_A+\Delta_B)/(\bar A+\bar B)$.
Then $\ell^2=(1-\theta)/(pq)$ and $\lambda=n\E U^2=npq\theta/(1-\theta)$.
The \emph{weighted} versions of \eqref{j:local:re:residual} prove
\begin{equation}\label{j:local:re:exports}
 n\theta\le\frac{Nr_I}{K},\quad
 nd_Y\le\frac{Nr_I}{6(1-\Theta)(A_l+B_l)},\quad
 \lambda\le\frac{Ns_+r_I}{(1-\Theta)K}.
\end{equation}
Indeed $nD\le Nd$ whereas $nb_I\ge Nb_I$.
These exports are derived from the weighted resource; multiplying an
unweighted variance ceiling by $N$ would not prove them.

For $j=3,4,5$ set
\[
 L_3(v)=R+2v,\quad L_4(v)=R^2+2Rv+3v^2,
 \quad L_5(v)=R^3+2R^2v+3Rv^2+4v^3.
\]
Put $v_+=\sqrt{(1-l)/l}$, $v_-=\sqrt{u/(1-u)}$ and
$k_j=[1-(1-\Theta)^{j/2}]/\Theta$, with $k_j=j/2$ at $\Theta=0$.
For $u<1/2$, valid Bernoulli moment ceilings are, with $s=l(1-l)$,
\[
 B_3^*=(1-2s)/\sqrt s,\quad B_4^*=1/s-3,
 \quad B_5^*=(1-4s+2s^2)/s^{3/2}.
\]
They decrease towards $p=1/2$. Convexity places $(1-\theta)^{j/2}$
below its endpoint chord on $[0,\Theta]$; positive-measure pricing gives
\begin{equation}\label{j:local:re:moments}
 M_j\le B_j^*+r_I\max\left\{0,
 \frac{L_j(v_+)-k_jB_j^*}{K_+},
 \frac{L_j(v_-)-k_jB_j^*}{K_-}\right\}.
\end{equation}
The same rule, using
$T_+\le2mv_+\Delta_B+L_3(v_+)\theta_+$ and its reflection, bounds
\begin{equation}\label{j:local:re:outerresource}
 W_o:=(\sigma-1)T_++(\sigma+1)T_-
 \le r_I\max\left\{\begin{gathered}
 \frac{(\sigma-1)v_+}{3(A_l+B_l)},
 \frac{(\sigma+1)v_-}{3(A_l+B_l)},\\
 \frac{(\sigma-1)L_3(v_+)}{K_+},
 \frac{(\sigma+1)L_3(v_-)}{K_-}\end{gathered}\right\}.
\end{equation}
Intersect every new bound with the previous one. The versions which
merge $\theta_+$ and $\theta_-$ are obtained by replacing both costs by
their maximum and both prices by their minimum; they need no new proof.

\paragraph{The enclosure theorem.}\label{j:local:re:operator}
On each cell, the formula atlas bounds the four coefficients in
\eqref{j:core:remainder}. The result, for each actual sign $s=\operatorname{sgn}H$,
is a finite family of quadratics
\begin{equation}\label{j:local:re:quadratics}
 \mathcal R\le\max_j(A_j+F_{j,s}|H|-Q_j^{\rm raw}H^2),
 \qquad Q_j^{\rm raw}>0.
\end{equation}
The maxima retain all score--moment corners. In the initial signed
cover, intervals contain both signs of $\beta,\gamma$; no orientation is
discarded. In the probability cover, the central integral bound is a
maximum of affine functions of $\rho$, hence convex. The other
$\rho$-dependence is affine or the maximum of convex quadratics. For
fixed $H$, the score dependence is convex. Endpoints are therefore taken
\emph{before} completing the square.

\begin{theorem}[Finite resource enclosure]\label{j:local:acquisition}
Consider a finite ordered schedule of the following operations:
\begin{enumerate}
\item Split a cell into an explicitly covered family of closed cells.
\item Intersect its enclosures with an inequality proved in the atlas.
\item Exclude a cell by an empty necessary interval or a strictly
negative nonnegative resource.
\item Apply a separately proved support implication,
checking all its hypotheses against this cell's already acquired data.
\item Export a raw, weighted, or signed phase bound by
\eqref{j:local:re:quadratics}, \eqref{j:local:budget}, and \eqref{j:local:re:exports}.
\end{enumerate}
If exact outward arithmetic checks every domain, corner and cover,
every output contains all attaining laws satisfying the input. The
conclusion holds for all integers $n\ge N$ wherever the formula's
monotonicity gates are checked. It includes zero conditional variance,
singleton intersections and arbitrarily small atoms.
\end{theorem}
\begin{proof}
Induct over the schedule. A split preserves the union, intersections use
necessary inequalities, and an excluded cell contains no feasible law.
Each callback consumes only earlier hypotheses for that same law.
For a scalar raw output, use
$fh-qh^2\le(f_+)^2/(4q)$ for $h\ge0$, $q>0$; using $f^2$ is a valid
weaker bound. For a resource output first add \eqref{j:local:re:extra} to each
curvature. A negative maximum excludes the cell; a nonnegative one may
be divided by $\kappa N$ to give its next defect parameter.
The phase bound follows from the positive root of
$A+fh-qh^2\ge0$, replacing $A$ by $A_+$ if necessary.
A strict output is exported only when the certificate retains a strict
upper gap, such as its strict Fourier or target enclosure; closed-domain
coverage does not change strictness. All constants are frozen before evaluating
corners. Thus no iteration
step relies on an unproved fixed point. The stated index uniformity is
that of the displayed finite formulas, not a fixed-law limit.
\end{proof}

For the Fourier cancellation estimate only, reflect $(X,H)$ if needed:
$I_0(-X,-H)=I_0(X,H)$. Both odd coefficients change sign, and
$\int e^{-x^2/2}|zx^2-cx^4|\,dx$ is unchanged by $(z,c)\mapsto(-z,-c)$.
Consequently a lower endpoint used there bounds $|\gamma|$, not necessarily
the original signed $\gamma$. The actual contact, gap and phase data keep
their original orientation. A useful additional signed lower bound is
\begin{equation}\label{j:local:re:gamma}
 \gamma\ge(1-\Theta)^{3/2}\frac{1-2u}{\sqrt{u(1-u)}}
       -\left(R+2\sqrt{\frac{u}{1-u}}\right)\Theta.
\end{equation}
To prove it expand $\E X^3=\E V^3+\E[W^2(X+2V)]$; the positive
residual is nonnegative and the negative residual has the displayed
lower bound. The binary term decreases in both $p<1/2$ and $\theta$.
Thus $\max(0,\text{any signed lower bound})$ is a valid lower bound for
$|\gamma|$ inside this auxiliary Fourier estimate.



\paragraph{Signed moments with a common first moment.}\label{j:local:atlas}\label{j:local:moments}
Put $a_\pm=(1\pm\beta)/2$, $b_\pm=(\rho\pm\gamma)/2$.
Cauchy gives $m^2\le a_+a_-$ and $b_\pm m\ge a_\pm^2$.
Consequently
\[
 \frac{(1+\beta)^{3/2}}{\sqrt{1-\beta}}-\rho
 \le\gamma\le
 \rho-\frac{(1-\beta)^{3/2}}{\sqrt{1+\beta}},\qquad
 m\ge\max\left\{\frac{1+\beta^2}{2\rho},
                  \frac{a_+^2}{b_+},\frac{a_-^2}{b_-}\right\}.
\]
These statements also hold with a zero atom. Projection of $X|X|$ onto
the span of $X$ and $|X|$ gives
\[
 \mu_4\ge\rho^2+\frac{(\gamma-\rho\beta)^2}{1-\beta^2}.
\]
For completeness, this follows by minimizing the nonnegative expectation
$E(X|X|-uX-v|X|)^2$ over $u,v$; its Gram matrix has diagonal entries
one and off-diagonal entry $\beta$, and is positive definite.

The signwise moment matrices below, valid also for different radii,
provide the bounded-support fourth- and fifth-moment contractions.
A second useful upper bound for $m$ follows by integrating
$(x-a)^2(R-x)\ge0$ for $x=|X|\in[0,R]$. For every fixed $a>0$,
\[
 2m\le\frac{R-\rho+2a+Ra^2}{a(2R+a)}.
\]
The cover uses only the eleven rational choices $a=2/10,\ldots,12/10$.

Write $\delta=\gamma-\rho\beta$. Cauchy and a square completion on
each sign give the signed constant bound
\begin{align*}
 \frac\rho2\mu_4-\frac14\eta_5
 &\le\frac{\rho^3}{4}
 -\frac{(\rho+\gamma)(\delta_+)^2}{8(1+\beta)^2}
 -\frac{(\rho-\gamma)((-\delta)_+)^2}{8(1-\beta)^2}.
\end{align*}
Indeed on either sign, write $a=\E X_s^2$, $b=\E X_s^3$,
$c=\E X_s^4$, $f=\E X_s^5$. Cauchy gives $c\ge b^2/a$,
$f\ge c^2/b$, and
$\rho c/2-f/4\le\rho^2b/4-(c-\rho b)^2/(4b)$.
Distance to the half-line $c\ge b^2/a$ proves the stated penalty
on summing both signs. Independently,
$\mu_4^2\le\rho\eta_5$ and the projection bound above give
\[
 \frac\rho2\mu_4-\frac14\eta_5
 \le\frac{\rho^3}{4}
 -\frac{(\gamma-\rho\beta)^4}{4\rho(1-\beta^2)^2}.
\]
The larger of the two nonnegative deductions is retained.

\paragraph{An actual central gap without assuming \texorpdfstring{$D<2$}{D below 2}.}
If zero belongs to the support, the signwise gap calculation gives $J\le\rho$.
Thus $D\ge\sqrt{10}\rho-\gamma\ge(\sqrt{10}-1)\rho$.
Whenever $D<(\sqrt{10}-1)\rho$, there is therefore a central gap
$(-A,B)$, $A,B>0$, with no zero atom. Put $p=P(X>0)$, $q=1-p$,
$T_+=b_+-mB^2$ and $T_-=b_--mA^2$. The exact gap decomposition is
\[
 D=mQ+(\sigma-1)T_++(\sigma+1)T_-+3(T-J_o),\qquad
 Q=(\sqrt{\sigma+1}A-\sqrt{\sigma-1}B)^2,
 \quad\sigma=\sqrt{10}.
\]
All terms are nonnegative. Hence $T\le D/(\sigma-1)$ and
\[
 (A+B)^2\ge A^2+B^2\ge\frac{\rho-D/(\sigma-1)}m.
\]
Moreover $A\le a_-/m$, $B\le a_+/m$ and $J\le\rho+2m$.
The last bound gives the compulsory defect
\[
 D\ge \sigma\rho-\gamma-6m.
\]
The inequalities $T-J_o\ge2BR_++2AR_-$ give the useful additional
endpoint constraint
\[
 D\ge(\sigma-1)b_++6a_+B-(\sigma+5)mB^2.
\]
If $D<(\sigma-1)b_+$, this downward quadratic forces $B$ to be at
least its positive root. The reflected constraint has $a_-,b_-,A$
and the coefficients $\sigma+1,\sigma+7$. In a closed moment box,
positive coefficients are rounded down, negative coefficients up, and
the lower square-root endpoint is used in the resulting root bound.

The actual opposite-sign event has probability $2pq$.
The first two moment constraints give lower bounds for both $p$ and $q$.
For independent copies one has exactly
\[
 E[(X-X')^2\mid XX'<0]
 =\frac{1+\beta(1-2p)+4m^2}{2pq}
 =\frac{1+4m^2+\beta}{2p}
  +\frac{1+4m^2-\beta}{2q}.
\]
For an upper bound, use the upper $m$ endpoint and both endpoints of
$p$ and $\beta$. The expression is convex in $p$ and affine in $\beta$.
This improves the former $2/P(XX'<0)$ bound inside the complete HIGH
estimate. Both copies still have the same actual marginal.

\paragraph{Preserving central and outer integral costs.}
Let
\[
 P_1(A,B)=B^3/3+AB^2/2-A^3/6,\quad
 P_2(A,B)=B^4/4+AB^3/3-A^4/12,\quad G=\rho P_1-P_2.
\]
Then $J_{1,c}=mP_1$, $J_{2,c}=mP_2$. Since
$\partial_BG=B(A+B)(\rho-B)$,
\[
 G(A,B;\rho)\le G(A,\rho;\rho)
 =\rho^4/12+\rho^3A/6-\rho A^3/6+A^4/12.
\]
The last expression increases with $A$ because its derivative is
$(\rho-A)^2(\rho+2A)/6$, and is convex in $\rho$ because its second
$\rho$ derivative is $\rho^2+\rho A$. Its closed-rectangle maximum is
therefore bounded using the upper $A$ endpoint and both $\rho$ endpoints.

Positive outer gaps satisfy
$(t-a)(\rho-t)\le(\rho-B)_+^2/4$. On negative outer gaps the positive
part is bounded by $(R-\rho)^2/4$ after extracting $|t|$.
Integration gives
\[
 3(\rho J_1-J_2)_o
 \le\frac38\{(\rho-B)_+^2T_++(R-\rho)^2T_-\}.
\]
The available weighted outer budget is $D-mQ$, not the entire $D$.
The cover intersects this budget with the individual bounds
$T_+\le b_{+,u}-m_lB_l^2$ and $T_-\le b_{-,u}-m_lA_l^2$.

For the signed phase term $U=J_1+(\mu_4-3)/6$, retain
\[
 -RT_-/2\le J_{1,o}\le RT_+/2.
\]
The central minimum of $P_1$ occurs at $B_l$ and an $A$ endpoint;
its maximum occurs at $B_u$ and $A$ clipped to the interval at $A=B_u$.
A nonnegative polynomial bound is multiplied by $m_u$ for an upper
bound and $m_l$ for a lower bound; a negative one uses the opposite
endpoint. For example, the binary endpoints $A=3/2$, $B=2/3$ make this
maximum negative; the sign switch is necessary even when $D<2$.
The fourth-moment sign intervals also sharpen
$K_1=E[X|X|^3]-3\beta-3\rho\gamma/2$.
The complete $K_1,U,\phi$ rectangles are used in the signed expression
$|CK_1+\phi U|$, without replacing it by two separate absolute values.


\paragraph{A complete coarse bound before a central gap is licensed.}
Always $-\mu_4/6\le J_1\le\mu_4/2$ and
$-3J_2-\eta_5/4\le0$. On a gap of length $h=b-a$, these follow from
\[
 \int_a^b(t-a)t\,dt+\tfrac12\int_a^bt^2dt=hb^2/2,\quad
 \tfrac32\int_a^bt^2dt-\int_a^b(t-a)t\,dt=h(2a+b)^2/6,
\]
and, on a negative gap $(-v,-u)$, from
$\tfrac13\int_u^v x^3dx-\int_u^v(v-x)x^2dx=u^3(v-u)/3$.
For $1\le r\le\rho$ and actual support radius $R\le6$, set
\[
 \Gamma(R,r)\ge\max_{r\le v\le R}
 \left\{-\frac{v^2}4+\frac{(R+r)v}3-\frac{Rr}2+
                      \frac{r^3(2R-r)}{12v^2}\right\}_+.
\]
Then
\[
 3(\rho J_1-J_2)\le\frac32\Gamma(R,r)(\rho-\gamma)
                              +\frac38\rho^2J.
\]
On negative gaps extract $x=|t|$ and discard the nonpositive part
$x<\rho$. Replacing $\rho$ by $r$ and the outer endpoint by $R$
gives $\int_r^v(R-x)x(x-r)\,dx$. Dividing this integral by $v^2$
gives exactly the displayed expression. Signwise Tonelli now charges
$\Gamma\E X_-^3=\Gamma(\rho-\gamma)/2$.
On positive gaps $(t-a)(\rho-t)\le\rho^2/4$, giving $\rho^2J/8$.
The implementation uses 64 closed $v$-intervals: on each, maximize the
concave quadratic at its clipped vertex $2(R+r)/3$ and the positive
reciprocal term at the left endpoint. This proves its interval bound.
The alternative bound $\Gamma\le5/2$ follows at $r=1,R=6$ by
$3v^4-28v^3+66v^2-11>0$ on $[1,6]$: the endpoint values are $30,205$,
and the interior critical values are $(85\mp35\sqrt5)/2>0$.
These monotone enclosing parameters cover all $r\ge1,R\le6$.
\paragraph{Moment and gap coefficients for a whole cell.}
\label{j:local:mom:conditional-boxes}

The following coefficient rules are collected here to make the finite
calculation transparent. Every rule is an inequality for an actual law;
there is no extremal-law assumption inside a cell.
Write $v=1-\theta$. For $p\in[l,u]\subset(0,1)$ set
\begin{gather*}
 s_0=\min\{l(1-l),u(1-u)\},\quad
 v_* =\max\{\sqrt{u/(1-u)},\sqrt{(1-l)/l}\},\\
 r_B=(1-2s_0)/\sqrt{s_0},\quad
 g_B=\max_{t=l,u}|1-2t|/\sqrt{t(1-t)},\quad b_B=\max_{t=l,u}|1-2t|,\\
 L_3=R+2v_*,\quad L_4=R^2+2Rv_*+3v_*^2,\\
 L_5=R^3+2R^2v_*+3Rv_*^2+4v_*^3,\quad
 e_B=r_B\frac{1-4s_0+2s_0^2}{s_0(1-2s_0)}.
\end{gather*}
The polynomials $T_3=x+2a$, $T_4=x^2+2ax+3a^2$ and
$T_5=x^3+2ax^2+3a^2x+4a^3$ are nonnegative. Their maxima at
$x=R$, $a=v_*$, in Lemma~\ref{j:local:re:pricing}, give, if $\theta\le\Theta<1$,
\begin{align}
 \rho&\le r_B+(L_3-r_B)\Theta,&
 |\gamma|&\le g_B+(L_3-g_B)\Theta,\notag\\
 |\beta|&\le b_B+(1-b_B)\Theta,&
 \mu_4&\le(1-\Theta)^2(1/s_0-3)+L_4\Theta,\notag\\
 \eta_5&\le(1-\Theta)^2e_B+L_5\Theta.&&\label{j:local:mom:boxes}
\end{align}
Replacing $\theta$ by $\Theta$ uses the explicitly checked inequalities
$L_3>r_B,g_B$, $L_4>2(1/s_0-3)$ and $L_5>2e_B$.
The exact binary factors are $v^{3/2},v,v^2,v^{5/2}$; the upper bounds
use $v^{3/2}\le v$ and $v^{5/2}\le v^2$. Thus cells crossing $1/2$ and
zero conditional variances are included. The same cubic expansion, in the
later domain $p\in[.37,.47]$, $A\ge.75$, $B\ge1.05$, gives
\begin{equation}\label{j:local:mom:rho-lower}
 \rho\ge v^{3/2}\rho_B(p)+3\min(A,B)\theta
       \ge\rho_B(p)+.45\theta,\qquad
 \rho_B(p)=\frac{1-2pq}{\sqrt{pq}}.
\end{equation}
Indeed $\rho_B<1.11$ there and $v^{3/2}\ge1-3\theta/2$.
Also $\beta\ge1-2p-2(1-p)\theta$.

For the signed coefficient
$K_1=\E[X|X|^3]-3\beta-3\rho\gamma/2$, define
\[
 K(p,v)=(1-2p)\{(3v^3/2-v^2)/(pq)-3v^3+v^2+3v\}.
\]
Direct substitution gives $K_1(V)=-K(p,v)$, with the coefficient $3$
unchanged. For $v\ge2/3$, $K$ decreases in $p$: both coefficients in
braces are nonnegative and $(1-2p)/(pq)$ decreases. Its derivative in $v$
has the sign of $1-2p$, since the derivative of the braces is at least
$9v^2-6v+3>0$. Hence both variance endpoints are retained. The error is
\begin{equation}\label{j:local:mom:K-error}
 |K_1(X)-K_1(V)|\le
 \Theta\{L_4+3+\tfrac32L_3(R_\rho+g_B)\}
\end{equation}
when $\rho\le R_\rho$. This follows from the conditional expansions and
$\rho\gamma-\rho_V\gamma_V=
\rho(\gamma-\gamma_V)+\gamma_V(\rho-\rho_V)$.

Here are the remaining signed charges. They also explain the two
alternative bounds retained in the initial, coarse moment cover:
\begin{align}
 -\mu_4/6&\le J_1\le\mu_4/2,&
 -3J_2-\eta_5/4&\le0,                                  \label{j:local:mom:gap-signs}\\
 \rho J_1-J_2&\le\tfrac54(\rho-\gamma)+\tfrac18\rho^2(\rho+1),
                                                              \label{j:local:mom:oldgap}\\
 -75D/104&\le J_{1,o}\le3D/2.                           \label{j:local:mom:outer}
\end{align}
The first two lines follow from the gap integral identities above,
using $\Gamma\le5/2$ and $J\le\rho+1$. For the outer $J_1$ bounds,
$J_{1,o}\ge-3T_-$, $J_{1,o}\le3J_o\le3(T_++T_-)$ and
\eqref{j:core:gap-defect} give the stated constants.

For the previously defined $P_1,P_2$, $J_{j,c}=mP_j(A,B)$. Both polynomials increase in $B$; their $A$ derivative changes from positive to negative
at $A=B$. A rectangle minimum is therefore at $B_l$ and an $A$ endpoint.
Multiply a negative minimum by an upper bound for $m$, and a nonnegative
minimum by a lower bound. For the upper central charge, maximizing at
$(\bar a,\bar b)$ if $p\le1/2$ gives
$J_{1,c}\le v^2(2-3p)/(6pq)$, a decreasing function of $p$ on that interval.
If $p\ge1/2$, maximize first at $A=B$ to get
$J_{1,c}\le(2/3)v^2q^2/p\le1/3$. Thus every cell may use
\begin{equation}\label{j:local:mom:central-upper}
 J_{1,c}\le\max\{1/3,(2-3l)/(6l(1-l))\}.
\end{equation}

For $|X|\le R$, $R\ge1$, define
\[
 k_4(R)=\frac{(R+1)^2}{R+1/2},\qquad
 k_5(R)=\frac{R^3+2R^2+3R+3/2}{R+1/2}.
\]
Then
\begin{equation}\label{j:local:mom:universal-moments}
 \mu_4\le\min\{R\rho,1+k_4(R)(\rho-1)\},\quad
 \eta_5\le\min\{R\mu_4,1+k_5(R)(\rho-1)\},\quad
 \mu_6\le R\eta_5.
\end{equation}
The first two bounds are the supported moment estimates proved in
Section~\ref{j:fourier:compensated}; the last is pointwise.
The following coarse cancellation estimates retain the radius-six majorants.
The last identity in~\eqref{j:local:re:pricing} gives
\begin{equation}\label{j:local:mom:K-cancellation}
 |K_1+3\rho\beta|\le(96-18\rho)(\rho-1)/13.
\end{equation}
Indeed $u(u+2-3r/2)/(u+1/2)$ is convex on $[0,6]$, its maximum is
$(96-18r)/13$, and its negative part is at most $31/40$, smaller than that
maximum for $1\le r\le37/20$.
The following bounds are used only for $|X|\le6$ and
$1\le\rho\le37/20<2$, with the signed moment hypotheses below.
The other bounds needed on this domain are
\begin{align}
 |K_1|&\le6\rho-3\rho^2/2+3|\beta|,\notag\\
 K_1&\le M+3|\beta|-3\rho g_-/2,&
 -K_1&\le3\rho+3|\beta|+(3\rho/2-3)g_-,\notag\\
 -K_q&\le(\rho/2-3/8)M-\rho/2+1/8+3|\beta|g/2,&
 -K_q&\le8\rho^3/27+(3|\beta|/2-1/2)\rho,               \label{j:local:mom:K-bounds}
\end{align}
where $\mu_4\le M$, $|\gamma|\le g$, $\gamma\ge g_-$, and
$K_q=\eta_5/4-\rho\mu_4/2-3\beta\gamma/2+\rho/2$.
For the first two lines use
$|\E[X|X|^3]-3\gamma|\le3\rho$, and note the negative coefficients
of $\gamma$. For the last line use
$\eta_5\ge\mu_4^{3/2}\ge3\mu_4/2-1/2$, or maximize
$\rho y^2/2-y^3/4$ over $y\ge0$.


\paragraph{Compatible weaker resource constants.}
The implementation also permits the weaker binary charge
$b_I=(616/25)(1-\Theta)^2\operatorname{dist}(p_*,I)^2$.
It follows from $2(6+2\sigma)>616/25$ and
$(1-\theta)^{3/2}\ge(1-\theta)^2$.
With $K=\min(K_+,K_-)$ and $v_* =\max(v_+,v_-)$, the merged charge
\eqref{j:local:re:residual} gives
\[
 T_++T_-\le T_*=(d-b_I)\max\{v_*/[3(A_l+B_l)],(R+2v_*)/K\},
 \qquad d_o=\min\{d,(\sigma+1)T_*\}.
\]
Thus $W_o\le d_o$ and, for $R\le6$,
\[
 -75d_o/104\le J_{1,o}\le3d_o/2,\qquad
 3(\rho J_{1,o}-J_{2,o})\le\frac38
 \max\{\rho^2/(\sigma-1),(R-\rho)^2/(\sigma+1)\}d_o.
\]
These are precisely the merged-cost versions of the sign-specific
resource rule; they do not introduce another available budget.

\paragraph{Conditional lower moments improve two signed charges.}

Let $s_u=\max_{p\in[l,u]}pq$ and $\theta\le\Theta$.
The unsigned cases of~\eqref{j:local:re:conditional} give
\begin{align}
 \mu_4&\ge M_l:=\max\{1,(1-\Theta)^2(1/s_u-3)\},\label{j:local:entry201:m4lower}\\
 \eta_5&\ge E_l:=
 (1-\Theta)^{5/2}\frac{1-4s_u+2s_u^2}{s_u^{3/2}}.\label{j:local:entry201:e5lower}
\end{align}
Both binary moment functions decrease in $s\in(0,1/4]$.
For the fifth this follows because the derivative has the sign of
$s^2+2s-3/2<0$; the fourth is immediate.
The fourth-moment lower bound sharpens the lower endpoint of
$J_1+(\mu_4-3)/6$; its explicit use appears in the cell map below.

The constant signed charge is
\[
 -K_q=\frac{\rho\mu_4}{2}-\frac{\eta_5}{4}
       +\frac32\beta\gamma-\frac{\rho}{2}.
\]
The old coherent estimate used $\eta_5\ge3\mu_4/2-1/2$ and
$\mu_4\le M$. A second estimate uses \eqref{j:local:entry201:e5lower}.
Their difference is independent of $\rho$. Consequently the improved
bound is exactly the old one minus
\begin{equation}\label{j:local:entry201:eta-improvement}
 \frac14\max\{0,E_l-3M/2+1/2\}.
\end{equation}
This is a constant shift on the cell and does not spoil convexity.
It does not assert the generally invalid implication
$\eta_5\ge3M/2-1/2$ from an upper bound on $\mu_4$.

\paragraph{Signed conditional Taylor cancellation.}

Let $K_1=\E[X|X|^3]-3\beta-\tfrac32\rho\gamma$.
On a sign class put $x=|X|$, $a=|V|$, $W=X-V$ and
$\varepsilon=1$ on the positive class, $-1$ on the negative class.
In the actual central-gap domain there is no atom at zero.
With $w_3=x+2a$, $w_4=x^2+2ax+3a^2$ and
$\gamma_V=\E V^3$, exact conditional expansion gives
\begin{equation}\label{j:local:entry201:Kexact}
 K_1(X)-K_1(V)
 =\E\!\left[W^2
 \left\{\varepsilon(w_4-3-\tfrac32\rho w_3)
                   -\tfrac32\gamma_Vw_3\right\}\right].
\end{equation}
Apply~\eqref{j:local:re:conditional} to $j=2,3,4$, with $h=1$
or $h=\varepsilon$, and use
$\rho\gamma-\rho_V\gamma_V
=\rho(\gamma-\gamma_V)+\gamma_V(\rho-\rho_V)$.

The positive coefficient in \eqref{j:local:entry201:Kexact} is
$F(x,a;t)=x^2+2ax+3a^2-3-\tfrac32t(x+2a)$
with $t=\rho+\gamma_V$; the negative coefficient is $-F$
with $t=\rho-\gamma_V$. Their ranges are bounded on complete
rectangles of actual conditional means and support magnitudes.
The positive-support magnitude is at most $R$.
The same formula permits separate actual bounds on the two signed support magnitudes.
On the late domain $p<1/2$, bounds for $\gamma_V$ are
\[
 \sqrt{(1-\Theta)^3(1-2u)^2/[u(1-u)]}
 \le\gamma_V\le(1-2l)/\sqrt{l(1-l)}.
\]
The positive mean lies in
$[\sqrt{(1-\Theta)(1-u)/u},\sqrt{(1-l)/l}]$, and the negative
magnitude mean lies in
$[\sqrt{(1-\Theta)l/(1-l)},\sqrt{u/(1-u)}]$.
Use the lower bound $\rho\ge\rho_B(u)$ only after checking
$3\min(A_l,B_l)\ge(3/2)\sup_I\rho_B$. Indeed conditional expansion
and $(1-\theta)^{3/2}\ge1-3\theta/2$ then imply
$\rho\ge(1-\theta)^{3/2}\rho_B(p)+3\min(A,B)\theta\ge\rho_B(p)$.
Otherwise retain the universal lower bound one.

Since $x+2a>0$, $F$ decreases with $t$.
For fixed $t$ its Hessian is
$\left(\begin{smallmatrix}2&2\\2&6\end{smallmatrix}\right)>0$.
Its rectangle maximum is at a corner. Its minimum is among the four
corners, the minimizer on each edge, and the stationary point
$(x,a)=(3t/8,3t/8)$ if that lies in the rectangle.
On fixed-$x$ edges use $a=t/2-x/3$, clipped to the edge; on
fixed-$a$ edges use $x=3t/4-a$, also clipped.
This is a complete finite extremum list.
Multiplication by $0\le\theta\le\Theta$ uses
$\Theta\min(0,L)$ and $\Theta\max(0,U)$ for the lower and upper
coefficient bounds. It intersects this new interval with the retained
absolute-error interval; neither need be discarded.

\paragraph{The complete central rectangle.}
The central part of $\rho J_1-J_2$ is exactly
$mG(A,B;\rho)$, where $G=\rho P_1-P_2$.
Let a proved probability cell give
$A\in[A_l,A_u]$, $B\in[B_l,B_u]$, $m\in[m_l,m_u]$.
Here
\[
 A_u=\sqrt{u/(1-u)},\quad B_u=\sqrt{(1-l)/l},\quad
 m_l=\sqrt{(1-\Theta)s_0},\quad
 m_u=\sqrt{\max_{p\in[l,u]}pq}.
\]

\begin{lemma}[Complete central rectangle]\label{j:local:entry201:central}
Write $G(A_l+(A_u-A_l)x,B_l+(B_u-B_l)y;\rho)
=\sum t_{kh}x^ky^h$. Define the 25 rational coefficients
\[
 b_{ij}(\rho)=\sum_{\substack{k\le i\\h\le j}}
 t_{kh}\frac{\binom{i}{k}}{\binom4k}
          \frac{\binom{j}{h}}{\binom4h},\qquad 0\le i,j\le4,
\]
and
\[
 \mathcal G(\rho)=\max_{i,j}
       \{m_l b_{ij}(\rho),\,m_u b_{ij}(\rho)\}.
\]
Then $\rho J_{1,c}-J_{2,c}\le\mathcal G(\rho)$.
Moreover $\mathcal G$ is convex in $\rho$.
\end{lemma}
\begin{proof}
Apply Lemma~\ref{j:cert:bernstein} separately to the
two coordinates after the affine substitution. The Bernstein basis is nonnegative and sums to one on the closed square,
so $G$ is a convex combination of its 25 coefficients.
For each coefficient retain both $m$ endpoints, or equivalently use
$m_u$ when it is nonnegative and $m_l$ otherwise. This gives a uniform
bound on the entire three-dimensional rectangle, including boundaries.
Each endpoint-multiplied coefficient is affine in $\rho$; their maximum
is convex. No interior stationary point is omitted.
\end{proof}



\paragraph{One interval-feedback theorem.}
\label{j:local:entry:feedback}\label{j:local:fourier:uniform}

The signed identity from the preceding section is
\begin{equation}\label{j:local:entry:remainder}
 \cR=C\{3\beta J/2+3\rho J_1-3J_2-K_q\}
 +H(CK_1+\varphi U)-H^2(C\rho+\gamma\varphi/2)+E_n,
\end{equation}
where $U=J_1+(\mu_4-3)/6$, $\varphi=\phi(H/\sqrt n)$.
For any one coherent low-frequency triple and complete high-frequency value $K_H$ at floor $N$,
\begin{equation}\label{j:local:entry:fourier-budget}
 |E_n|\le K_F+s_F|H|,\quad
 K_F=c_\pi(L_0+K_H+L_2/\sqrt N),\quad s_F=c_\pi L_H,\quad c_\pi\ge2/\pi.
\end{equation}
The implementation uses $16/25$ before the signed phase and $212/333$
afterward; both are valid since $\pi>333/106$.

\begin{lemma}[Continuous cells and feedback]\label{j:local:entry:cell-theorem}\label{j:local:entry:corners}
Suppose bounds for one actual cell give
\[
 r_0\le\rho\le r_1,\quad0<c_0\le C\le c_1,\quad
 |\beta|\le b,\quad |\gamma|\le g,\quad \mu_4\le M,\quad D\le d,
\]
$K_1\in[k_l,k_u]$, $U\in[u_l,u_u]$, and
$0\le\varphi_-\le\varphi\le2/5$. Freeze the shared outer charge $d_o$,
lower fifth moment $E_l$, upper $s_u=\max_{p\in[l,u]}p(1-p)$, and
central envelope $\mathcal G$ from the preceding bounds.
Define
\begin{align*}
 A(r)&=\tfrac32 b(r+2\sqrt{s_u})+3\mathcal G(r)
 +\tfrac38\max\{r^2/(\sigma-1),(R-r)^2/(\sigma+1)\}d_o\\
 &\quad +(r/2-3/8)M-r/2+1/8+\tfrac32bg
 -\tfrac14\max\{0,E_l-3M/2+1/2\},\\
 B(c)&=\max_{k\in\{k_l,k_u\},\,u\in\{u_l,u_u\},\,
                      \psi\in\{\varphi_-,2/5\}}|ck+\psi u|,\\
 Q(c,r)&=cr+\tfrac12\varphi_-(3.16r-3-d)
       -\tfrac12(2/5-\varphi_-)(d-.16)_+.
\end{align*}
If all four corner values of $Q$ are positive, then, for every real $H$,
\begin{align}
 \cR&\le\max_{(c,r)\in\{c_0,c_1\}\times\{r_0,r_1\}}
       \{K_F+cA(r)+(B(c)+s_F)|H|-Q(c,r)H^2\}\notag\\
 &\le\max_{c,r\text{ corners}}
       \{K_F+cA(r)+(B(c)+s_F)^2/(4Q(c,r))\}.           \label{j:local:entry:four-corners}
\end{align}
For the following feedback, assume in addition that the cell contains a
ratio extremizer with $e=C_n-\CE\ge0$, $n\ge N$, and that
$0<T<.4N$ is a strict upper bound for the final maximum in
\eqref{j:local:entry:four-corners}. Then $T<nC_n\rho$, since
$C_n\rho\ge\CE>.4$, and the weight-removal step in
Corollary~\ref{j:core:license} gives
\begin{equation}\label{j:local:entry:feedback-values}
 D<\frac{15000T}{997N},\quad
 \varphi>\frac{997}{2500}\left(1-\frac{5T}{2N}\right),\quad
 C<\CE+V/N\quad\text{if also }\cR/\rho<V.
\end{equation}
Only positive density floors are used. These statements remain valid at
$H=0$, on closed cell edges, and for zero conditional variance.
\end{lemma}
\begin{proof}
The constant coefficient follows from the shared outer budget,
Lemma~\ref{j:local:entry201:central}, and~\eqref{j:local:mom:K-bounds}, using
\[
 J\le\rho+2ABm\le\rho+2(1-\theta)^{3/2}\sqrt{pq}
       \le\rho+2\sqrt{s_u}.
\]
This follows from $A\le\bar a$, $B\le\bar b$, $m^2=(1-\theta)pq$;
it is used only in the charge $3|\beta|J/2$. Since $\gamma\ge3.16\rho-3-d$, its negative
part is at most $(d-.16)_+$. Using the upper density only for that negative
part gives $C\rho+\gamma\varphi/2\ge Q$. Absolute value is separately
convex in each coordinate of $ck+\psi u$, proving the formula for $B$.
For fixed $r$ and $|H|$, the resulting upper bound is convex in $c$.
Take those endpoints \emph{before} completing the square. At either
endpoint $c>0$, it is convex in $r$, since $\mathcal G$ and the maximum of the two outer quadratics are
convex, every other term of $A$ is affine, and $Q$ is affine in $r$. This proves the first line of~\eqref{j:local:entry:four-corners};
$fh-qh^2\le f^2/(4q)$ proves the second. Bilinearity of $Q$ proves
its positivity throughout the rectangle from its corners.

\end{proof}

\paragraph{Signed phase curvature in the resource schedule.}
\label{j:local:re:signed-phase-repair}
The auxiliary Fourier orientation does not change the signed third moment
in the actual transfer identity. The paired phase estimate must therefore
retain its sign. On a fixed resource cell put
\[
 L=\frac{997}{2500}\left(1-\frac{5r}{2N}\right),\qquad
 g_0(\rho)=3.16\rho-3-d,\qquad U=.4,
\]
where $r,d$ are the preceding raw and defect upper bounds, and the
checker requires $0\le L\le U$. The actual density lies in $[L,U]$
and the actual signed moment satisfies $\gamma\ge g_0(\rho)$. Consequently
\begin{equation}\label{j:local:re:sign-aware-price}
 C\rho+\tfrac12\gamma\phi
 \ge q_{\rm s}(C,\rho):=
 C\rho+\tfrac12\min\{L g_0(\rho),U g_0(\rho)\}.
\end{equation}
This follows by first multiplying the lower bound on $\gamma$ by the
nonnegative actual density, and then minimizing over its two endpoints.
The simpler expression $C\rho+Lg_0(\rho)/2$ is licensed if $g_0\ge0$,
or if nonnegativity of the actual $\gamma$ has been separately proved.
It is not licensed merely by reflecting the auxiliary Fourier integral.

The corrected price preserves the four-corner reduction. For fixed $C$,
$q_{\rm s}$ is concave in $\rho$, being affine plus the minimum of two
affine functions. For fixed $\rho$ it is affine in $C$. In the phase
envelope $A+f|H|-q_{\rm s}H^2$, the remaining terms are separately
convex in $C,\rho$, by the moment and coefficient bounds already proved.
Thus the entire envelope is separately convex, and two successive
one-dimensional endpoint bounds place its maximum at the four corners.
The same argument applies after adding the affine conserved-resource
curvature $e_\rho$ and subtracting a barrier price $K\rho$.
If the four corrected divisors are positive, concavity makes them
positive throughout the rectangle, and square completion gives the
valid raw and resource bounds
\[
 \max_{C,\rho\text{ corners}}\left(A+\frac{f^2}{4q_{\rm s}}\right),
 \qquad
 \max_{C,\rho\text{ corners}}
       \left(A+\frac{f^2}{4(q_{\rm s}+e_\rho)}-K\rho\right).
\]
Any threshold-root export likewise uses the corrected corner curvature.

For any retained paired calculation using the nonnegative-moment
specialization, the conditional cubic bound \eqref{j:local:re:gamma}
is applied in the actual contact orientation. Its probability interval,
support radius and variance ceiling must all have been acquired for that
same marginal before the call. The selected domain checks give a lower
bound greater than $.096$, using the already acquired common radius
$R=3.84$. This licenses the specialization independently of the auxiliary
Fourier orientation.

There is also a convenient cell-constant correction. Since $\rho\ge1$,
$(-g_0(\rho))_+\le(d-.16)_+$, so the weaker price
\[
 C\rho+\tfrac12L g_0(\rho)
       -\tfrac12(U-L)(d-.16)_+
\]
is valid. Both forms keep the raw remainder separate from the
barrier-subtracted resource budget.

The broad preliminary and complementary resource passes are replayed with
\eqref{j:local:re:sign-aware-price}; their subsequent consumers retain the same
fixed target tables only when the complete corrected replay verifies every
exported resource. A diagnostic equality of two corner maxima alone does
not replace this check, since threshold bounds also consume the phase
coefficients. Later uses of the simpler formula require their own explicit
nonnegative signed-moment or small-defect license.




\paragraph{Using a floor calculation at a later index.}
\label{j:local:transfer}
A floor calculation encloses one fixed attaining ratio marginal with
$C\ge\CE$; it does not assert monotonicity of $C_n$. For actual
$n\ge N$, every saved pass
preserves the two simultaneous statements
\begin{equation}\label{j:local:ext:floor-invariant}
 \mathcal R_n\le r,\qquad nD\le Nd.
\end{equation}
In particular $D\le d$ may be used in index-free moment and gap
inequalities. The stronger weighted statement is retained whenever a
cost is proportional to the sample size.

Here is why each operation in the selected schedules preserves this
invariant. Keep the actual marginal's moments, support, signed phase
and admissible frequency windows fixed while increasing $n$.
All numerical gates mentioned below are checked before the operation.

\emph{Fourier bounds.} Positive sums and products of nonnegative
decreasing factors remain decreasing. For the complete HIGH bound
of Proposition~\ref{j:fourier:high}, these factors are
\[
 \sqrt n e^{-cn}\ (2Nc>1),\quad n^{-1},\quad
 e^{-knb^2}/\sqrt n,\quad
 r(x)=\frac{x+2}{x^2+4x+2},\quad x\text{ increasing in }n.
\]
Indeed $r'=-(x^2+4x+6)/(x^2+4x+2)^2<0$.
The comparator $\sqrt n e^{-n/8}\{2r(n/8)+g_c/3\}$ decreases
for $N\ge5$. The improved first piece has the same factors with
its separately checked positive rate and slope. For the optional zero
piece, the derivative of $\sqrt{x}e^{-x}r(x)$ has the sign of
\[
 4-10x-22x^2-13x^3-2x^4.
\]
This polynomial decreases for $x>0$ and equals $-11/128$ at
$x=1/4$, so the native gate $Nkb^2\ge1/4$ suffices. Every use
retains the event/window conditions $0<a\le1$, $m_e>0$,
$0<d_0\le12$, $1\le M\le111/10$, $0<w\le1/4$ and positive
sine slope. Taking the minimum of complete admissible windows is
valid even when the minimizing window changes with $n$.

For LOW, the nonconstant factors are
\[
 (n/(n-1))^{3/2},\quad (n/(n-1))^{3/2}/(n-1),\quad
 a_n^{-j},\quad b_n^{-j},\quad n^{-1/2},
\]
where $a_n=(n-1)/(2n)$, $b_n=(45-M)(1-2/n)/96$ and $j>0$.
They decrease on the stated positive domains. This covers the full
error $e/(2a_n^2)+B^2/(2b_n^3)$ and $L_2/\sqrt n$; the paired
moment-rectangle hypotheses still have to hold.

\emph{Density and score updates.} Under the raw sign and positivity
license $r<.4N$, the invariant \eqref{j:local:ext:floor-invariant} gives
\[
 \phi(H/\sqrt n)\ge\phi_{0,l}(1-5r/(2n)).
\]
Its lower bound increases with $n$, whereas
$r/n-\phi_{0,l}(1-5r/(2n))D_l/6$ decreases for $r,D_l\ge0$.
Both may therefore be evaluated at $N$. The general curvature
correction has derivative $(g_0+t)/2\ge0$ with respect to this
density floor, where $g_0=(79/25)\rho-3-d$, $t=(d-4/25)_+$:
if $d\le4/25$, use $g_0\ge4/25-d$; otherwise
$g_0+t=(79/25)(\rho-1)$. The signed curvature
\[
 C\rho+\tfrac12\min\{\phi_lg_0,(2/5)g_0\}
\]
also increases with $\phi_l$ when $g_0\ge0$ and is constant
when $g_0<0$. Thus both signs are retained. Positive added
curvatures and all divisors are checked before square completion.

For $C\ge\CE$ and $\CE<40974/10^5$, the compulsory charge is
\[
 n(C-\CE)\rho\ge
 N\max\{0,C-40974/10^5\}\rho_l.
\]
Split the score cover at $40974/10^5$; the charge is affine on
each side. The coherent corner calculation then exports
$nD\le B/\kappa$, hence $d=B/(\kappa N)$, for
$\kappa=997/15000$. A negative budget excludes the cell before
division. Positive $K/n$ calculations separately retain their native
$K\rho$ charge; two lower bounds on the same excess are not added.
Closed covers, upward rounding and intersection with an earlier
bound now preserve \eqref{j:local:ext:floor-invariant} by finite
induction, without a fixed-point or limiting argument.

\emph{Weighted physical resources.} If an acquired inequality is
\[
 D\ge b+g d_{\rm in}+(N/n)h_N\lambda_n,
 \qquad \lambda_n=ns,\quad b,g,h_N\ge0,
\]
then the invariant gives
$(n/N)(b+g d_{\rm in})+h_N\lambda_n\le d$, hence
$b+g d_{\rm in}+h_N\lambda_n\le d$.
This retains the original joint intercepts. The actual $n$ is still
used in every binomial probability. The transfers used here are
$N=25$ to $26\le n\le30$ and $N=31$ to $32\le n\le34$;
each actual index has its own verified support calculation.
Index-free support refinements use only already acquired data. For
example $6mAB\ge P>0$ gives $(A+B)^2\ge2P/(3m_u)$ by
$4AB\le(A+B)^2$; it does not establish an initial central gap.

\paragraph{Phase and variance exports from the same budget.}
Suppose one signed phase family gives, for $h=|H|\ge0$,
\[
 0\le D\le c_n\max_j(A_j+F_jh-Q_jh^2),\qquad c_n>0,
 \quad Q_j>0.
\]
Some corner is nonnegative at the actual $h$. For that corner
$\Delta_j=F_j^2+4Q_jA_j\ge0$ and
$h\le(F_j+\sqrt{\Delta_j})/(2Q_j)$.
If $u_j\ge0$ is a rational upper enclosure of $\sqrt{\Delta_j}$,
the complete family therefore gives
\[
 h\le\max_j r_j,\qquad
 r_j=\begin{cases}0,&\Delta_j<0,\\
 \max\{0,(F_j+u_j)/(2Q_j)\},&\Delta_j\ge0.
 \end{cases}
\]
This includes double roots and $h=0$. Keep separate maxima for
the two phase signs and intersect each with its previous cap.
At floor31 the normalization is $15000/(997\cdot31)$; its
actual-index version $c_n=15000/(997n)>0$ leaves this argument
unchanged. Positivity of every $Q_j$ is required.

Finally suppose bounded conditional supports imply $s\le B\delta$.
Writing $O_0,O_1$ for their outward offsets, take
$B=\max(O_0,O_1)\max(p_u,1-p_l)>0$.
Combining this with the original joint
plane $\delta/d_0+\lambda_n/L\le1$ and $\lambda_n=ns$ gives
\begin{gather*}
 \lambda_n\le\{L^{-1}+(nBd_0)^{-1}\}^{-1},\\
 \theta_n=\frac{\lambda_n}{npq+\lambda_n}
 \le\frac1{1+n v_l/L+v_l/(Bd_0)}
 \le\frac1{1+N v_l/L+v_l/(Bd_0)},
\end{gather*}
where $v_l=\min_{p\in[p_l,p_u]}pq>0$ and $d_0,L>0$.
Upward rounding and intersection with the previous variance bound
preserve this conclusion. A downward support-scale enclosure enlarges
$O_0,O_1$, so also weakens the bound safely. The refinement concerns
physical $\theta_n$; it does not identify $\lambda_n$ with a
floor's diagnostic $\lambda_N$ or replace the plane's intercept $L$.
All 193 saved first-band refinements at floor 25 are checked at each
actual index 26--30.


Here is the version of the moment matrices that permits different
positive and negative support radii. On either sign write $Y=|X|$,
$0\le Y\le r$, and let $p_s,m,a,b,c_4,c_5$ be its unnormalized
moments of orders zero through five. Positivity of the weights
$r-Y$, $Y(r-Y)$, and $Y^2(r-Y)$ gives
\[
 \begin{pmatrix}rp_s-m&rm-a\\rm-a&ra-b\end{pmatrix},\quad
 \begin{pmatrix}rm-a&ra-b\\ra-b&rb-c_4\end{pmatrix},\quad
 \begin{pmatrix}ra-b&rb-c_4\\rb-c_4&rc_4-c_5\end{pmatrix}
 \succeq0.
\]
For positive denominators their determinant inequalities imply
\begin{align*}
 p_s&\ge m/r+(rm-a)^2/[r(ra-b)],\\
 c_4&\le rb-(ra-b)^2/(rm-a),\\
 c_5&\le rc_4-(rb-c_4)^2/(ra-b).
\end{align*}
A zero diagonal entry forces the corresponding off-diagonal entry to
vanish. One may check that condition or discard the nonnegative
correction, but one cannot divide by zero. In interval form a correction
is bounded below by the square of a positive-part lower numerator
divided by an upper positive denominator. The sign radius $r$ is used
only for its own class. Summing the first inequality over the two signs
and using $p_++p_-\le1$ can improve the common upper mean.

One further bound used at small indices strengthens the signed constant
coefficient. Set $\delta=\gamma-\rho\beta$ and $Y=|X|$. Then
\[
 V:=\E\{Y(Y^2-\rho Y)^2\}
   =\E|X|^5-2\rho\mu_4+\rho^3,
 \qquad
 \delta=\E\{Y\operatorname{sgn}(X)(Y^2-\rho Y)\}.
\]
Here $m>0$ by centering and unit variance.
Cauchy's inequality for the nonnegative measure $Y\,dP$ gives
$\delta^2\le2mV$. Consequently
\begin{equation}\label{j:local:ext:weighted-cubic-loss}
 \frac\rho2\mu_4-\frac14\E|X|^5
 \le\frac{\rho^3}{4}-\frac{(\gamma-\rho\beta)^2}{8m}.
\end{equation}
This deduction can be combined with the two other deductions in
Section~\ref{j:local:atlas} by taking their maximum. Their sum would spend
the same nonnegative quantity more than once.



\paragraph{A universal lower bound for each sign probability.}
\label{j:local:ext:sign-mass}
Six probability panels at index six are excluded before the count
calculation. The exclusion is the following elementary moment inequality;
it does not require support, a central gap, or absence of a zero atom.

\begin{lemma}[Sign mass and third moment]\label{j:local:sign-mass}
If $\E X=0$, $\E X^2=1$, and $\rho=\E|X|^3<\infty$, then
\[
 \rho^2\Pp(X>0),\quad \rho^2\Pp(X<0)
 \ \ge\ \frac{8(7\sqrt7-17)}{27}
 \ >\ \frac{11261}{25000}.
\]
\end{lemma}
\begin{proof}
Write $p=\Pp(X>0)$, $a=\E X_+^2$, and
$m=\E X_+=\E X_->0$. Centering and unit variance imply
$0<a<1$. Cauchy's inequality, first with factors
$X_+^{1/2}$ and $X_+^{3/2}$ and then with their negative-sign
counterparts, gives
\[
 a^2\le m\E X_+^3,\qquad (1-a)^2\le m\E X_-^3.
\]
Also $m^2\le pa$. Consequently
\[
 p\rho^2\ge f(a):=\frac{(a^2+(1-a)^2)^2}{a}.
\]
The derivative has the sign of $6a^2-2a-1$, because
\[
 f'(a)=\frac{(2a^2-2a+1)(6a^2-2a-1)}{a^2}.
\]
Thus the unique minimum on $(0,1)$ is at
$a=(1+\sqrt7)/6$, where
$f(a)=8(7\sqrt7-17)/27$. The strict rational comparison follows
from $2645751^2<7\cdot10^{12}$ and
$8(7\cdot2645751/10^6-17)/27>11261/25000$.
Apply the same argument to $-X$ for the other sign probability.
\end{proof}



\paragraph{Scalar consequences of the common inward budget.}
\label{j:local:ext:inward-scalar}
Some finite count domains use a stronger absolute-third-moment floor or
noise ceiling before their five-dimensional subdivision. The following
calculation supplies those scalar restrictions. It uses the sign
conditioning of \eqref{j:core:conditional-moments}.

Reflect only for the scalar calculation, so that
$0<p\le1/2$, $q=1-p$, $v=pq$. This preserves $\rho$, total
conditional variance $s$ and total inward displacement. The count
probability, its phase and its original orientation are not reflected.
Suppose the inward endpoints obey $a_0+a_1\le d$ with $0<d\le p$.
The conditional third-moment inequalities and Cauchy's inequality give
\begin{equation}\label{j:local:ext:inward-I}
 \rho\ge I(p,s,d):=
 \frac{v(1-2v)+(3p-d)s+s^2/(qd)}{(v+s)^{3/2}}.
\end{equation}
Indeed, writing $b=qs_0$, $c=ps_1$ and $b+c=s$, the absolute-cubic
numerator is at least
\[
 v(1-2v)+3pb+3qc-a_0b-a_1c
                       +\frac{b^2}{qa_0}+\frac{c^2}{pa_1}.
\]
Its linear part is at least $(3p-d)s$, while its quadratic part is
at least $(b+c)^2/(q(a_0+a_1))\ge s^2/(qd)$.
Vanishing inward endpoints force the corresponding centered variance
to vanish; those terms are handled before division. If $d=0$, then
$s=0$ directly.

The derivative has the exact numerator
\begin{equation}\label{j:local:ext:inward-I-derivative}
 (v+s)^{5/2}I_s
 =v(6p-3p^2-d-3/2)
   +\left\{\frac{2p}{d}-\frac{3p-d}{2}\right\}s
   +\frac{s^2}{2qd}.
\end{equation}
Consequently, on $p\in[l,u]\subset(0,1/2]$, the two gates
$d\le l$ and $f:=6l-3l^2-d-3/2\ge0$ imply
\[
 \rho\ge\rho_B(p):=\frac{1-2p(1-p)}{\sqrt{p(1-p)}}
                  \ge\rho_B(u).
\]
The last function is decreasing on this half-interval. If a previously
acquired ceiling is $\rho\le R$, a point with
\[
 R^2p(1-p)-(1-2p(1-p))^2<0
\]
is impossible. This polynomial increases on $0<p\le1/2$, because
its derivative is $(1-2p)(R^2+4-8p(1-p))\ge0$.
Exact bisection may therefore raise the closed lower probability
endpoint to a known infeasible point; retaining that endpoint is a
conservative enlargement.

When $f>0$ and $s\le S$, formula~\eqref{j:local:ext:inward-I-derivative}
also gives
\[
 I_s\ge
 \frac{l(1-l)f}{[u(1-u)+S]^{5/2}}.
\]
Integration from zero supplies a new variance ceiling. Directed square
roots and outward division make every finite update valid separately.

A useful weaker floor survives failure of the monotonicity gate.
Assume still $0<d\le l$ and set
\[
 a=(3/2-6l+3l^2+d)_+,
 \qquad B=2l/d-3u/2\ge5/4.
\]
The only possible decrease of $I$ is bounded below by
$-(va-Bs)_+/v^{5/2}$. Integrating this negative triangle gives
\begin{equation}\label{j:local:ext:inward-deficit}
 \rho\ge\rho_B(u)-\frac{a^2}{2B\sqrt{l(1-l)}}.
\end{equation}
In particular, failure of the monotonicity gate must not be replaced
by the unsupported assertion $\rho\ge\rho_B(p)$.

There is a second variance contraction which works with either sign
of $f$. The global Taylor upper bound for $(v+s)^{3/2}$, combined
with \eqref{j:local:ext:inward-I}, gives the necessary inequality
\begin{equation}\label{j:local:ext:scalar-noise-quadratic}
 Gs^2+fs\le\epsilon\,[u(1-u)+S]^{3/2},
\end{equation}
where
\[
 \epsilon=(R-\rho_B(u))_+,\qquad
 G=\frac1{(1-l)d}-\frac{3(1-2l(1-l))}{8l(1-l)}
       \ge\frac5{8l(1-l)}+\frac34>0.
\]
To see this, subtract $\rho_B(p)(v+s)^{3/2}$ from the lower
numerator in \eqref{j:local:ext:inward-I}, use the quadratic Taylor upper
bound on the subtracted term, and then use $p\ge l$ and
$\rho_B(p)\ge\rho_B(u)$. The positive part in $\epsilon$ is
necessary when $f<0$. The larger real root of
\eqref{j:local:ext:scalar-noise-quadratic}, enclosed from above, is a valid
new upper bound for $s$, even for negative $f$.

Finally suppose the original joint resource has the stronger form
\[
 d_{\rm actual}\le d_0(1-s/S_0),\qquad S_0=\Lambda/n.
\]
To test a proposed upper endpoint $a\in(0,S_0)$, assume $s\ge a$.
Then $d_{\rm actual}\le d_0(1-a/S_0)$, so the preceding quadratic
calculation can be repeated with this smaller inward bound. If its
resulting upper bound is strictly less than $a$, the whole region
$s\ge a$ is impossible. Otherwise the probe proves no lower bound
on $s$. A finite sequence of such hypothetical tests is sound
without any convergence assertion. Throughout it, the original
$S_0$ remains in the joint budget; a separately improved noise
ceiling clips that budget and does not replace its intercept.


\subsection{Paired Fourier cancellation}
\label{j:local:paired}
The triangle LOW bound separates the fourth- and fifth-order errors.
Near a two-point law their leading imaginary parts partly cancel.
We retain that cancellation in one Gaussian norm; all other Taylor and
power remainders remain those of Proposition~\ref{j:fourier:low}.

\begin{lemma}[The cancellation norm and its finite tail bounds]
\label{j:local:entry189:mills}
Put $I=\sqrt{\pi/2}$ and
\[
 J(z,c)=\int_0^\infty e^{-x^2/2}|zx^2-cx^4|\,dx.
\]
The function $J$ is convex in $(z,c)$ and invariant under simultaneous
reflection. Its maximum on a rectangle is therefore bounded by its
largest corner value. For $zc\le0$ it equals $I(|z|+3|c|)$.
For $z,c>0$, put $x=z/c$ and define
\[
 d_{k+1}=1,\qquad d_j=1+\frac{j}{xd_{j+1}}\quad(j=k,\ldots,1),
 \qquad m_k(x)=d_1^{-1}.
\]
Then $J$ is bounded by each of
\begin{align}
 U(z,c)&=I|z-3c|+6\sqrt{zc}\,e^{-x/2},\notag\\
 V(z,c)&=I(z-3c)+\sqrt{zc}\,e^{-x/2}
       \left(6+\frac{2(3-x)}x m_k(x)\right),
 \label{j:local:entry189:mills-bound}
\end{align}
where $k=11$ for $x\ge3$ and $k=12$ otherwise. The parenthesized factor
is positive. For negative $z,c$, apply these formulas to $|z|,|c|$.
\end{lemma}
\begin{proof}
Convexity follows by integrating the absolute value of a linear function;
successive endpoint bounds give the rectangle assertion. For $zc>0$,
reflection and integration by parts, split at $r=\sqrt{z/c}$, give
\[
 J=I(z-3c)+6cr e^{-r^2/2}
       +2(3c-z)\int_r^\infty e^{-t^2/2}\,dt.
\]
Dropping a nonpositive last term, or bounding its integral by $I$ when
positive, proves $U$. To retain the signed tail, set
$I_j(r)=\int_0^\infty t^j e^{-rt-t^2/2}\,dt$ and
$q_j=I_j/I_{j-1}>0$. Integration by parts yields
\[
 I_0=\frac1{r+q_1},\qquad q_j=\frac j{r+q_{j+1}},\qquad
 I_0=e^{r^2/2}\int_r^\infty e^{-t^2/2}\,dt.
\]
Truncating at $q_{k+1}=0$ reverses the inequality at every reciprocal:
odd truncations are lower bounds, even ones upper bounds. At $r=\sqrt x$
their values multiplied by $r$ are exactly $m_k(x)$ above. The tail's
multiplier $3c-z$ selects the stated parity and proves $V$, without
requiring convergence of an infinite fraction. For $x<3$ the
parenthesized factor exceeds $6$; for $x\ge3$, $0<m_{11}(x)<1$ makes it
greater than $6-2(x-3)/x>4$.
\end{proof}

For outward evaluation use an upper enclosure for $I$ in $U$ and in
the case $zc\le0$. In $V$, use a lower enclosure when $z-3c<0$
and an upper one otherwise; the remaining factor is positive. Thus
square roots and $e^{-x/2}$ may be bounded upward. Let $J_u$ be the
largest such corner bound, taking the minimum of $U,V$ when applicable.

\paragraph{A single transfer formula.}
Retain the hypotheses and notation of Proposition~\ref{j:fourier:low}:
$N>2$, $1\le\mu_4\le M\le11.1$, $|\gamma|\le g\le1.85$, and actual
ceilings $|a|\le a_*$, $|z|\le Z$, $\mu_6\le M_6$, where
\[
 a=(3-\mu_4)/24,\qquad z=(10\gamma-\E X^5)/120,
 \qquad c=a\gamma/6.
\]
Any rectangle containing the actual pair $(z,c)$ may be used above.
For $F_N\ge(N/(N-1))^{3/2}$ and $I_u\ge I$, a valid LOW triple is
\begin{align}\label{j:local:fourier:paired}
 L_0&=F_NJ_u+\frac{F_NI_u}{N-1}
       (35a_*g^3/48+35Zg^2/24)+251g^3/17280,\\
 L_H&=F_NI_u\{a_*+5Zg/[2(N-1)]\}+251g^2/4800,\notag\\
 L_2&=E/(2a_0^2)+B^2/(2b_0^3),\notag
\end{align}
with the same $E=M_6/720+g^2/72+(1/2+g^2/144)^3/5$,
$B=a_*+Z/2+E/4$, $a_0=(N-1)/(2N)$, and any
$0<b_0\le(45-M)(1-2/N)/96$.

Indeed, write $b=n-1$ and $v=b\gamma t^3/6$. The leading imaginary
part in the common Taylor proof, after division by $t^3$ and rotation
by $Ht$, is bounded by
\[
 |zt^2-atv|+a_*t|\sin v-v|+Zt^2|\cos v-1|
       +|H|t\{a_*t+Zt^2|\sin v|\}.
\]
Multiply by $ne^{-bt^2/2}$ and integrate. The first term contributes
$nJ(z,c)/b^{3/2}$; the others use
$|\sin v-v|\le|v|^3/6$, $|\cos v-1|\le v^2/2$ and
$|\sin v|\le|v|$. The recurrence
$\int_0^\infty x^{2j}e^{-x^2/2}\,dx=(2j-1)!!I$
gives respectively the coefficients $35/48,35/24,1,5/2$ in
\eqref{j:local:fourier:paired}. The common power/Taylor errors and
the signed Gaussian comparator give exactly the unchanged $L_2$ and
the two $251$ terms. All inverse-index factors decrease for $n\ge N$;
the conclusion is uniform when the law and threshold vary with $n$.

\paragraph{One moment envelope for the coefficient rectangle.}
The broad triangle choice is already proved in
Section~\ref{j:fourier:pearson-envelope}. With its exact radius-six
and actual-radius quantities $Z_0,Z_r$, it uses
\begin{equation}\label{j:local:fourier:z0}
 120Z=\min\{E_5+10g,Z_0,Z_r\},\qquad
 a_*=\max(2,M-3)/24,\quad M_6=RE_5.
\end{equation}
Here $E_5\ge\E|X|^5$ and $|X|\le R\le6$ are previously acquired.
The radius-six derivative test and the actual-radius clipped vertex
are exactly those in that proof; no moment ceiling is substituted
from an unrelated marginal. A positive rectangle gives the stronger
paired choice.

\begin{proposition}[A positive cancellation rectangle]
\label{j:local:paired-rectangle}
Suppose $\E X=0$, $\E X^2=1$, $|X|\le R$, $1\le\mu_4\le M<3$,
$0\le\gamma_l\le|\gamma|\le g$, $g^2<2$, $R>2g$, and
$M-1-\gamma_l^2\ge0$. After reflecting both odd Fourier coefficients
if necessary, set
\begin{gather*}
 t_* =\min\{g,(10-2M)/(2(R-g))\},\\
 120Z=\min\{E_5+10g,R(M-1)+(10-2M)t_*-(R-g)t_*^2\},\\
 z_l=\max\{-Z,[8\gamma_l-\gamma_l^3
        -(R+2\gamma_l)(M-1-\gamma_l^2)]/120\},\quad z_u=Z,\\
 a_*=(2-\gamma_l^2)/24,\qquad
 c_l=(3-M)\gamma_l/144,\quad c_u=(2-\gamma_l^2)g/144,\\
 M_6=\min\{RE_5,1+3g^2+g^4+
 (R^2+2Rg+3g^2+2)(M-1-\gamma_l^2)\}.
\end{gather*}
These are valid coefficient bounds and $[z_l,z_u]\times[c_l,c_u]$
contains $(z,c)$. If $3g^2\le2$, one may further decrease $c_u$ to
$\min\{c_u,(2g-g^3)/144\}$. Use the damping constant
$b_0=(45-M)(1-2/N)/96$; for $M\le2$ the weaker
$43(1-2/N)/96$ is also valid.
\end{proposition}
\begin{proof}
For $Q=X^2-\gamma X-1$, the common Pearson identities and one further
multiplication give
\begin{align}\label{j:local:fourier:pearson}
 \E Q^2&=\mu_4-1-\gamma^2,&
 \E X^5&=\gamma^3+2\gamma+\E[Q^2(X+2\gamma)],\notag\\
 \mu_6&=1+3\gamma^2+\gamma^4+
       \E[Q^2\{(X+\gamma)^2+2\gamma^2+2\}].
\end{align}
Thus $\mu_4\ge1+\gamma^2$, proving $a_*,c_l,c_u$; the optional
upper bound uses increasing $2t-t^3$ on $[0,g]$.
The positive-kernel estimates of Section~\ref{j:fourier:pearson-envelope}
give $-Q_R(t)\le120z\le P_R(t)$ for $t=\gamma\ge0$.
Since $M<3$, their absolute upper envelope is
$R(M-1)+(10-2M)t-(R-g)t^2$, maximized at $t_*$.
Moreover $(-Q_R)'(t)=10-2M+2Rt+3t^2>0$, proving $z_l$.
In the last identity, bound each nonnegative factor upward using
$t\le g$ and $\E Q^2\le M-1-\gamma_l^2$ to obtain $M_6$.
The transfer's domain follows from $M<3<11.1$, $g<\sqrt2<1.85$.
\end{proof}

Finally a probability row must supply a lower bound for the actual
$\rho$, not automatically its Bernoulli counterpart. If
$p\in[l,u]\subset(0,1/2)$, $\theta\le\Theta<1$, and acquired inner
endpoints are at least $a\ge0$, the conditional cubic identity gives
$\rho\ge(1-\theta)^{3/2}\rho_B(p)+3a\theta$.
For outward $r_-\le\rho_B(u)$, $r_+\ge\rho_B(l)$ this implies
\[
 \rho\ge r_*:=\max\{1,r_--(3r_+/2-3a)_+\Theta\},
\]
because $\rho_B$ decreases on $(0,1/2)$ and
$(1-\theta)^{3/2}\ge1-3\theta/2$.
The same $r_*$ is retained in all Fourier corners, signed phase prices,
conserved budgets and the signed fourth-moment errors
\eqref{j:local:entry201:Kexact}; the binary-$K$ alternative still requires
$1-\Theta>2/3$. Reflection here concerns only the auxiliary norm and
does not change the actual contact orientation.

\subsection{The acquired law and the comparison strategy}
\label{dc:entry}

The support, moment and Fourier arguments preceding this section place a
possible maximizing law near two points. Their output includes a bounded
conditional error and a common budget for that error. We now compare its
sum directly with the Bernoulli sum having the same two class probabilities.
The purpose of the comparison is to rule out the maximizing law, not to
describe every law with the same first few moments.

There are two useful features of the argument. First, one coupling controls
all changes in the Bernoulli count through a function of a single conditional
sum. Secondly, the expectation of a quartic upper bound for that function
is determined by four moments of the same conditional laws. A finite
endpoint argument then leaves polynomial inequalities on the original
probability and phase cells. No small positive noise cutoff is required.

Write $q=1-p$, $v=pq$, and
\[
 Y=B+U,\qquad \Pp(B=1)=p,\qquad \E(U\mid B)=0,
 \qquad 0<p<\tfrac12.
\]
The actual conditional laws $U_0,U_1$ satisfy
\begin{equation}\label{dc:support}
 -A_0\le U_0\le a_0,\quad -a_1\le U_1\le A_1,\qquad
 a_i\ge0,\quad a_0+a_1\le d_0<p,\quad A_0,A_1\ge d_0.
\end{equation}
The outward bounds need not be smaller than $1/2$. Put
\[
 s_i=\E U_i^2,\quad s=qs_0+ps_1,\quad \lambda=ns,
 \quad w=v+s,\quad \ell=w^{-1/2},\quad \theta=s/w.
\]
At the actual threshold choose a nearest integer $k$:
\[
 x=np+H/\ell=k+u,\qquad |u|\le\tfrac12,\qquad \delta=k-np.
\]
Both choices at a tie are allowed. The acquired threshold bounds are
checked before rounding and give $0<k<n$ on all applications below.
In particular, truncating a proposed count range to $[0,n]$ cannot remove
a possible threshold. All intervals used in this section are closed.

\paragraph{Inputs from acquisition.}
For the uniform entry hypotheses, the initial contact radius is $3.84$.
The central preparation and contact strip produce 45 rows on
$[.38,.47]$ with $A_0=A_1=.55$. The complementary preparation gives
39 left rows on $[.302,.38]$, with radius $3$ in the standardized
coordinate $X=\ell(Y-p)$. On a cell with $w\le w_+$ this gives the safe
outward bounds
\[
 A_0\ge3\sqrt{w_+}-p_-,\qquad
 A_1\ge3\sqrt{w_+}-(1-p_+).
\]
These support bounds and the common resource budget are the inputs to
the comparison. The acquisition theorem also supplies the complete
initialization and complementary probability exclusions.

\subsubsection{One shared price for inward displacement and noise}
\label{dc:budget}
The joint defect inequality already proved by acquisition is
\begin{align}
 D\ge{}&(1-\theta)^{3/2}D_B(p)
 +6\ell(1-\theta)\{d(1-d)+a_0a_1\}+b\theta,
 \qquad d=a_0+a_1,\label{dc:joint}\\
 D_B(p)={}&\frac{(6+2\sqrt{10})(p-p_*)^2}{\sqrt{pq}},
 \qquad p_*=(4-\sqrt{10})/2.\notag
\end{align}
Here $b$ is a lower bound supplied by the same signed row; one may take
the minimum of $3(\sqrt{10}+1)B_-$ and
$3(\sqrt{10}+3)A_-$, where $A_-,B_-$ are its inner endpoint bounds.
Let $P_0\le\inf_I D_B$, $\theta\le\Theta$, $d\le d_0$, and $v\le v_+$.
Choose positive rational lower bounds
\[
 \iota\le6(1-\Theta)\sqrt{(1-\Theta)/v_+}(1-d_0),\qquad
 \beta\le b-\tfrac32P_0.
\]
Then
\begin{equation}\label{dc:linear-budget}
                 D\ge P_0+\iota d+\beta\theta.
\end{equation}
Indeed $(1-\theta)^{3/2}\ge1-3\theta/2$ and
$d(1-d)+a_0a_1\ge(1-d_0)d$. This also explains why the two
contributions must be retained together.

For an actual index $n$, the acquired row gives
$nD\le\mathcal B$, where $\mathcal B$ is the minimum of its weighted
defect budget and its signed contact budget. If the acquisition was performed
at floor $N$, the former is $ND_{\rm row}$, not $nD_{\rm row}$.
The latter may be refined on a phase cell: enclose
$H=(k+u-np)/\sqrt{v+\lambda/n}$ with its sign, maximize each of the four
concave resource quadratics at its clipped vertex, and intersect with the
scalar resource bound. This operation uses the original common budget.

Set $E=\mathcal B-nP_0$. If $E<0$, the cell is excluded. Centering and
\eqref{dc:support} give $s_0\le A_0a_0$, $s_1\le A_1a_1$ and hence
$s\le Md$, for any $M\ge\max(qA_0,pA_1)$ on the cell.
Given an already proved ceiling $\lambda\le\Lambda$, one obtains
\begin{equation}\label{dc:noise-cap}
 \lambda\le
 \frac{E}{\iota/M+\beta/(v_++\Lambda/n)},\qquad
 d\le\frac{E}{n\iota}.
\end{equation}
These are simultaneous consequences of \eqref{dc:linear-budget}; no
resource has been counted twice. Each refinement starts from a proved
ceiling, so a finite number of refinements needs no convergence theorem.
All quotients use proved positive denominators. When $E=0$, positive noise
is impossible. When $P_0=0$, the inequality $nP_0\le\mathcal B$ gives no
upper index; the cell containing $p_*$ must not be discarded.

\subsection{One exact identity and one coupled payoff}
\label{dc:identity-section}
Let
\[
 \pi_j=\binom njp^jq^{n-j},\quad F_B(k)=\Pp(\operatorname{Bin}(n,p)\le k),
 \quad \rho_B=(1-2v)/\sqrt v,\quad
 \rho=\E|Y-p|^3/w^{3/2}.
\]
Use the positive closed-event score
\[
 C=\frac{\sqrt n}{\rho}
 \{\Pp(\textstyle\sum Y_i\le k+u)-\Phi((\delta+u)/\sqrt{nw})\},
 \qquad C_B=\frac{\sqrt n}{\rho_B}\{F_B(k)-\Phi(\delta/\sqrt{nv})\}.
\]
Reflection is applied before this construction when needed; the negative
branch at an atom uses the corresponding left limit, not the closed value
on the wrong side of that atom. Put $b_k=\sqrt n\pi_k$ and
$Q=\Pp(\sum Y_i\le k+u)-F_B(k)$.
For a fixed positive target $T$, direct subtraction gives
\begin{align}\label{dc:target-identity}
 \rho(T-C)={}&\rho_B(T-C_B)+\mathcal N-b_kQ/\pi_k+T(\rho-\rho_B),\\
 \mathcal N={}&\sqrt n\{\Phi((\delta+u)/\sqrt{nw})-
                                      \Phi(\delta/\sqrt{nv})\}.\notag
\end{align}
Here $T=\CE$ for the sharp comparisons and
$T=\CE+741/(4000\cdot39)$ for the separate positive-$39$ prerequisite.
The multiplier is this fixed target, even if the hypothetical score $C$
is larger. This distinction matters when the cubic correction is negative.

\begin{lemma}[One conditional sum controls all incoming counts]
\label{dc:coupling}
Suppose \eqref{dc:support} holds, $c\le1-d_0$ is positive, and
$R_0\ge\max_{j>k}\pi_j/\pi_k$ over legal counts. Define
\[
 \psi(z)=-\mathbf1_{\{z>0\}}+R_0\max\{0,\lfloor-z/c\rfloor\}.
\]
If $W_k$ sums $k$ independent copies of $U_1$ and $n-k$ copies of $U_0$,
then $Q/\pi_k\le\E\psi(W_k-u)$.
\end{lemma}
\begin{proof}
Take independent pairs $(U_{0,i},U_{1,i})$ and use them to couple every
conditional sum $W_j=\sum_{i\le j}U_{1,i}+\sum_{i>j}U_{0,i}$.
Successive values of $j+W_j$ increase by
$1+U_{1,j+1}-U_{0,j+1}\ge c$. Thus the successful counts form an
initial interval. If the cutoff is below $k$, its lost binomial mass is
at least $\pi_k$, and this event is exactly $W_k>u$.
If there are $r\ge1$ incoming counts above $k$, then
$W_k+cr\le u$, so $r\le\lfloor(u-W_k)/c\rfloor$.
Their combined mass is at most $rR_0\pi_k$.
Taking expectations gives the claim. Counts below zero or above $n$ have
mass zero. At $W_k=u$ the outgoing event is absent; at
$W_k=u-jc$ the incoming floor includes the equality. This proves the
closed-event statement, including every collision.
\end{proof}


A quartic $P(z)=\sum_{j=0}^4p_jz^j$ satisfying $P\ge\psi$ on the
whole real line therefore bounds all counts at once. Its coefficients
may depend on the closed cell but not on the unknown moments within it.
The next reduction evaluates its expectation from the same two conditional
laws that appear in \eqref{dc:target-identity}.

\subsection{Why only twenty-four or forty-eight moment tests remain}
\label{dc:corners}
For positive conditional variances write
$\gamma_i=\E U_i^3/s_i$; a zero-variance term vanishes and its
$\gamma_i$ can be assigned any allowed value. Support-weighted averages
give $-A_0\le\gamma_0\le a_0$ and $-a_1\le\gamma_1\le A_1$.
Moreover, integrating the nonnegative polynomials
$(U_0+A_0)(a_0-U_0)U_0^2$ and
$(U_1+a_1)(A_1-U_1)U_1^2$ gives
\begin{align}\label{dc:fourth}
 \E U_0^4&\le(a_0-A_0)\E U_0^3+A_0a_0s_0,\\
 \E U_1^4&\le(A_1-a_1)\E U_1^3+A_1a_1s_1.\notag
\end{align}
For $W=W_k$, independence and centering imply
\[
 t:=\E W^2=(n-k)s_0+ks_1,\quad
 \E W^3=(n-k)s_0\gamma_0+ks_1\gamma_1,
\]
\[
 \E W^4=(n-k)\E U_0^4+k\E U_1^4+3t^2
                         -3\{(n-k)s_0^2+ks_1^2\}.
\]
We may drop the last nonpositive term in an upper bound for $\E P(W-u)$
provided $p_4\ge0$.

The support signs also give the exact absolute-cube identity
\begin{equation}\label{dc:cube-exact}
 \E|Y-p|^3=v(1-2v)+sG,\qquad
 G=(1-\alpha)(3p-\gamma_0)+\alpha(3q+\gamma_1),
 \quad \alpha=ps_1/s.
\end{equation}
For $s>0$, $0\le\alpha\le1$. Taylor's theorem for $(1+z)^{3/2}$,
whose third derivative is negative for $z\ge0$, yields
\begin{equation}\label{dc:cube-lower}
 T(\rho-\rho_B)\ge
 \frac{T}{w^{3/2}}\frac\lambda n
 \left(G-\frac32(1-2v)-\frac{3(1-2v)\lambda}{8nv}\right).
\end{equation}
This formula retains the sign of the bracket.

Suppose, throughout a closed cell, rational constants satisfy
\begin{gather}\label{dc:enclosures}
 \rho_B(T-C_B)\ge g,\quad
 \mathcal N\ge f_0u+v_0\lambda,\quad
 0<b_-\le b_k\le b_+,\quad
 0<\chi_-\le T/w^{3/2}\le \chi_+,\\
 \lambda\in[L,U],\quad 0\le L\le U,\quad
 0<p_-\le p\le p_+<\tfrac12,\qquad
 -\tfrac12\le u_-\le u\le u_+\le\tfrac12.\notag
\end{gather}
The phase coefficient $f_0$ is a lower derivative bound on $u\ge0$ and
an upper derivative bound on $u\le0$. The scalar bounds are fixed on the
cell before eliminating the moments.

\begin{theorem}[Exact finite-cell comparison]\label{dc:finite-test}
\label{j:local:direct-finite}
Suppose \eqref{dc:support} and \eqref{dc:enclosures} hold, and choose
$c,R_0$ as in Lemma~\ref{dc:coupling}. Suppose $P\ge\psi$ globally,
with $p_4\ge0$. For each row of the following table,
set $t=V\lambda$, $M_0=1$, $M_1=0$, $M_2=t$, $M_3=\gamma t$,
and $M_4=J_4t+3t^2$:
\begin{center}
\begin{tabular}{ccccc}\toprule
$\alpha$&$\gamma$&$J_4$&$V$&$G$\\\midrule
0&$-A_0,\ 0,\ d_0$&$\gamma^2$&$(n-k)/(nq)$&$3p-\gamma$\\
1&$A_1,\ 0,\ -d_0$&$\gamma^2$&$k/(np)$&$3q+\gamma$\\\bottomrule
\end{tabular}
\end{center}
For both $\lambda\in\{L,U\}$, both $b\in\{b_-,b_+\}$ and both
$\chi\in\{\chi_-,\chi_+\}$, require positivity throughout the $p,u$ rectangle of
\begin{align}\label{dc:gate}
 g+f_0u+v_0\lambda
 &+\frac{\chi\lambda}{n}
 \left(G-\frac32(1-2v)-\frac{3(1-2v)\lambda}{8nv}\right)\\
 &-b\sum_{j=0}^4p_j\sum_{i=0}^j\binom ji(-u)^{j-i}M_i.\notag
\end{align}
These 48 tests imply $C<T$ for every actual law on the cell. If the cubic
bracket is independently proved nonnegative everywhere, only $\chi=\chi_-$ is
needed, leaving 24 tests.
\end{theorem}
\begin{proof}
First assume $\lambda>0$, so the variance proportion $\alpha$ is defined.
Combine \eqref{dc:target-identity}, Lemma~\ref{dc:coupling},
\eqref{dc:fourth} and \eqref{dc:cube-lower}.
An unknown positive multiplier in $[\chi_-,\chi_+]$ times a signed bracket is
at least the lesser endpoint product; this explains the two $\chi$ tests.
Fix either endpoint $\chi$. Parametrize
$s_0=\lambda(1-\alpha)/(nq)$ and $s_1=\lambda\alpha/(np)$.
For fixed $a_i,\gamma_i$ the resulting lower bound is affine in
$\alpha$ except for $-3bp_4t^2$, and is therefore concave in $\alpha$.
Its minimum is bounded by the endpoints $\alpha=0,1$.

At $\alpha=0$ only class $0$ remains. The expression is affine in
$\gamma_0$, so use $-A_0$ or $a_0$. The fourth coefficient at the first
endpoint is $(a_0-A_0)(-A_0)+A_0a_0=A_0^2$; at the second it is $a_0^2$.
In the latter case the only quadratic dependence on $a_0$ is
$-bp_4t a_0^2$, so endpoints $a_0=0,d_0$ suffice. Class $1$ gives the
other three entries in the same way. This is an extremal bound on the
moment expression; it does not assert $\E U_i^4\le\gamma_i^2s_i$ for an
actual conditional law.

At each corner the coefficient of $\lambda^2$ in the lower bound is
\[
 -\frac{3\chi(1-2v)}{8n^2v}-3bp_4V^2\le0.
\]
Thus the two noise endpoints suffice. The expression is affine in $b$,
which explains its two endpoints. These arguments enlarge necessary
moment domains and never choose different laws for different counts.
At $\lambda=0$ the continuous formulas give $M_i=0$ for $i\ge1$ and
require no division by a vanishing variance. Positivity of the final
lower bound gives $\rho(T-C)>0$, and $\rho>0$ finishes the proof.
\end{proof}

\subsection{What a quartic certificate proves}
\label{dc:checking}
The preceding theorem separates two claims. The polynomial must dominate
a discontinuous payoff on the whole real line, and its expectation must
leave a positive score deficit on a closed parameter cell. Both claims
are exact polynomial inequalities; neither follows from checking samples.

For the translated variable $z=W_k-u$, a sufficient global test is, for
an integer $J\ge1$,
\begin{align}\label{dc:global-tests}
 P(z)+1&\ge0&& (z\ge0),\\
 P(z)-jR_0&\ge0&& (-(j+1)c\le z\le-jc),\quad 0\le j<J,\notag\\
 P(z)+R_0z/c&\ge0&& (z\le-Jc).\notag
\end{align}
At $z=-jc$, the interval to its left supplies the higher step value;
at $z=0$, the $j=0$ interval supplies $P(0)\ge0$.
Thus closed adjacent tests include every event plane. The last ray
dominates the floor by its linear upper envelope. The selected finite
objects use $J=8$; this number is a format choice, not a probabilistic
assumption.

A compact interval is mapped affinely to $[0,1]$. On a ray, substitute
$z=z_0\pm t/(1-t)$ and multiply by $(1-t)^4$; the multiplier is positive
for $0\le t<1$, and the endpoint $t=1$ controls the leading term.
These transformations leave quartics with rational coefficients.
One acceptance method expands each target in the Bernstein basis and
checks nonnegative coefficients on a complete subdivision. A second,
more compact method accepts a strict target when its two endpoint values
are positive and an exact Sturm sequence counts no root in $(0,1)$.
The sequence starts with $F,F'$ and successive negative Euclidean
remainders. At a zero of an interior sequence member, its neighbors have
opposite signs, so the sign-variation count is unchanged; at a simple
root of $F$, it decreases by one. This proves the root count.
Zeros are omitted when counting endpoint sign variations. A positive
constant target is accepted directly, without a root sequence.
The selected strict reader rejects endpoint zeros and nonconstant gcds,
so repeated-root cases require no convention. Such rejections can make
discovery fail but cannot create a false positivity claim.

This strict reader is useful on a finite-index cell because its binary
reserve can pay for a small positive constant added to the quartic.
The uniform tail is different: its loss must remain proportional to
arbitrarily small noise, so its explicit polynomial has $P(0)=0$ exactly.
Its tangencies, including the repeated root of $Q_0+1$ below, are proved
by factorization and direct inequalities. They are not submitted to the
strict Sturm reader, which would reject those endpoint or gcd cases.

For the score tests, multiplying \eqref{dc:gate} by $(npq)^2>0$ gives a
polynomial of degree at most $6$ in $p$ and $4$ in $u$. The shift expansion
in that equation is part of the fixed checker; it is not an additional
moment estimate. On a probability/phase rectangle the tensor Bernstein
basis is nonnegative and sums to one. Positive coefficients therefore
prove the strict score inequality everywhere on the rectangle.
The change of basis is checked by exact expansion, and a cover tree must
contain both closed children at every split. A missing root, missing
child or unproved resource exclusion rejects.

Discovery may solve floating-point linear programs for the five quartic
coefficients. It may adaptively split a cell when a proposed polynomial
fails. Acceptance reads only rational coefficients and the full cover,
reconstructs the global payoff and all score targets, and checks their
inequalities exactly. There is no assertion that this search always
terminates. A particular completely accepted forest is sufficient.

\subsection{Scalar enclosures and the finite Bernoulli reserve}
\label{dc:normalization}
Here are explicit choices for the scalar quantities in
\eqref{dc:enclosures}. They also identify what the numerical calculations
must establish independently of the quartic coefficients.

Fix an actual integer $n$ and a sign-homogeneous phase cell. Let
$p\in[l,h]\subset(0,1/2)$, $\lambda\le U$, $v_-=l(1-l)$,
$v_+=h(1-h)$ and $w_+=v_++U/n$. Write
$\delta_-=k-nh$, $\delta_+=k-nl$, and let $M$ bound
$|\delta+t|$ for all $t$ between $0$ and every phase in the cell.
Choose rational $0<\phi_-\le\phi(0)\le\phi_+$ and set
\[
 a_\phi=\phi_-\left(1-\frac{M^2}{2nv_-}\right)>0.
\]
Since $e^{-y}\ge1-y$ for $y\ge0$, valid phase coefficients are
\begin{equation}\label{dc:phase-coeff}
 f_0\le a_\phi/\sqrt{w_+}\quad(u\ge0),\qquad
 f_0\ge\phi_+/\sqrt{v_-}\quad(u\le0).
\end{equation}
These bounds integrate the derivative of the first normal increment
in \eqref{dc:target-identity}. A lower coefficient for a negative phase
would reverse the required inequality.

For the remaining normal variance increment, differentiation in $s$ gives
\[
 \frac{d}{ds}\left[\sqrt n\Phi\left(\frac\delta{\sqrt{n(v+s)}}\right)\right]
       =-\frac{\delta\phi(\delta/\sqrt{n(v+s)})}{2(v+s)^{3/2}}.
\]
If $\delta_+<0$, use a positive lower bound on the density and on
$-\delta$ and the upper denominator $2w_+^{3/2}$; otherwise use
$-\phi_+\max(0,\delta_+)/(2v_-^{3/2})$.
Divide this bound by $n$ to obtain $v_0$ in
$\mathcal N\ge f_0u+v_0\lambda$.
Likewise, directed square roots give
\[
 \chi_-\le T_-/w_+^{3/2},\qquad \chi_+\ge T_+/v_-^{3/2},
 \qquad T_-\le T\le T_+.
\]
Both endpoints remain in the tests when the cubic bracket has either sign.

For fixed $n,k$, the mass $\pi_k(p)$ increases up to $p=k/n$ and decreases
afterward. Its minimum on $[l,h]$ occurs at an endpoint and its maximum at
an endpoint or the clipped mode. Exact rational evaluation at these points
and directed bounds on $\sqrt n$ give $b_-,b_+$.
Every ratio $\pi_j/\pi_k$, $j>k$, increases in $p$, so evaluating all legal
ratios at $h$ gives $R_0$. These facts avoid any asymptotic normal
approximation in the finite certificates.

\subsubsection{A quantitative binary gap on the whole probability interval}
\label{dc:binary}
The fixed target may be above the binary score even if the noise
perturbation itself is unfavorable. We therefore retain a quantitative
reserve rather than assume that each noisy law improves on its Bernoulli
counterpart. For a rational $p=a/M$, each binomial probability is an exact
integer divided by $M^n$; the integer mass recurrence and cumulative sums
give both $F_B(k)$ and $F_B(k-1)$.

The normal CDF can also be enclosed by rational arithmetic. For rational
$z\ge0$, integrate the Taylor expansion of $e^{-t^2/2}$:
\[
 \Phi(z)=\tfrac12+\phi(0)
 \sum_{j=0}^{m}\frac{(-1)^jz^{2j+1}}{2^jj!(2j+1)}+R_m,
 \qquad
 |R_m|\le\phi(0)\frac{z^{2m+3}}{2^{m+1}(m+1)!(2m+3)}.
\]
The remainder follows from the signed Taylor remainder for $e^{-x}$ on
$x\ge0$. Symmetry handles $z<0$, monotonicity handles an interval for $z$,
and directed square roots enclose $(k-np)/\sqrt{npq}$.
An alternating arctangent expansion in Machin's identity supplies rational
bounds for $\pi$ and hence $\phi(0)$ and $\CE$.

Two elementary continuous extensions of endpoint computations are used.
First, on $p\in[l,h]\subset(0,1/2)$, both $F_B(k)$ and
$z_k(p)=(k-np)/\sqrt{npq}$ decrease. Thus
\[
 F_B(k;p)-\Phi(z_k(p))\le F_B(k;l)-\Phi(z_k(h)),
\]
and the reflected jump satisfies
$\Phi(z_k(p))-F_B(k-1;p)\le\Phi(z_k(l))-F_B(k-1;h)$.
Multiply positive upper bounds by the increasing factor
$\sqrt{npq}/(1-2pq)$ at $h$. A finite regular closed grid supplies a
bound on the entire source interval, including both sides of every atom.

For narrower gaps, use a second-derivative chord bound. If $|f''|\le M$
on $[l,h]$, then
\begin{equation}\label{dc:chord}
 f(p)\le\max\{f(l),f(h)\}+M(h-l)^2/8.
\end{equation}
Indeed adding $Mp^2/2$ makes $f$ convex, and comparison with its chord
gives the stronger error $M(p-l)(h-p)/2$.
For the binomial part,
\[
 F_B(j;p)''=n\binom{n-1}j p^j(1-p)^{n-1-j}
                 \frac{(n-1)p-j}{p(1-p)}.
\]
The density factor attains its maximum at its clipped mode. For the normal
part, direct differentiation gives
\[
 z'=-\frac{k+(n-2k)p}{2\sqrt n\,v^{3/2}},\quad
 z''=-\frac{n-2k}{2\sqrt n\,v^{3/2}}
       +\frac{3[k+(n-2k)p](1-2p)}{4\sqrt n\,v^{5/2}}.
\]
Thus $|(\Phi\circ z)''|\le\phi_+(|z''|+|z||z'|^2)$, with every
quantity bounded on the full interval. This proves \eqref{dc:chord}'s
input without a sampled derivative estimate.

If the resulting cap gives $\CE-C_B\ge\epsilon$, then
$T-C_B\ge\epsilon+(T-\CE)$. To obtain $g$ in
\eqref{dc:enclosures}, multiply by a lower bound for $\rho_B$ when this
reserve is nonnegative and by an upper bound when it is negative.
A negative conservative reserve is allowed: the final quartic test,
not the sign of this intermediate estimate, decides the cell.
\subsection{An explicit quartic for the central tail}
\label{dc:tail}
For very large indices the finite binary reserve tends to zero. A bound
proportional to the noise is more useful there. The central rows allow one
explicit polynomial for every index $n\ge1000$, in place of a growing
family of finite-index certificates.

The polynomial
\begin{equation}\label{dc:q0}
                  Q_0(z)=-z-\tfrac12z^2+\tfrac12z^4
\end{equation}
has a simple purpose. Its linear term matches the Gaussian phase after
translation by $u$. Its negative quadratic term creates a positive
variance contribution in the score deficit. The quartic term controls
large negative errors and hence all incoming count steps. On the outgoing
side it has the exact tangency
\[
 Q_0(z)+1=(z-1)^2(z^2/2+z+1)\ge0.
\]
The finite-index correction changes only the scale and the quadratic
coefficient, while preserving its exact value zero at the event plane.

\begin{lemma}[A matched-slope global majorant]\label{dc:explicit-majorant}
For $a,c>0$, $R_0\ge0$, let
\[
 t_0=ac,\qquad
 \eta=\max\{0,(1-t_0^2)/2\}
               +\frac{\pos{R_0/t_0-1}}{t_0},\qquad
 P(z)=Q_0(az)+\eta(az)^2.
\]
Then $P(z)\ge-\mathbf1_{\{z>0\}}+R_0\max(0,\lfloor-z/c\rfloor)$
for every real $z$, and $P(0)=0$.
\end{lemma}
\begin{proof}
The factorization above handles $z>0$. For $-c<z\le0$,
$Q_0(-y)=y(1-y/2+y^3/2)\ge0$; split $0\le y\le1$ and $y\ge1$.
For $z\le-c$, put $y=-az\ge t_0$. Then
\[
 P(z)-R_0y/t_0
   =y(1-R_0/t_0)+y^2\{\eta+(y^2-1)/2\}\ge0.
\]
The first part of $\eta$ makes the brace nonnegative; its second part
pays a negative linear term since $y\ge t_0$. This dominates the linear
upper envelope of the incoming floor and includes its knots.
\end{proof}

For a symmetric outward support $|U_i|\le A$, write
$t=\E W_k^2$. Centering and independence give
$|\E W_k^3|\le At$ and $\E W_k^4\le A^2t+3t^2$.
Define the normalized phase increment
\[
 f(u)=\frac{\Phi((\delta+u)/\sqrt{nw})-\Phi(\delta/\sqrt{nw})}{\pi_k}.
\]
If $a_\ell\le f'\le a_h$, use $a=a_\ell$ for $u\ge0$ and
$a=a_h$ for $u\le0$. Then $f(u)\ge au$ on either phase half.
The linear term in $\E P(W_k-u)$ is exactly $au$, so cancellation gives
\begin{align}\label{dc:tail-loss}
 f(u)-\E P(W_k-u)\ge\mathcal L(t,h)
 :={}&a^2(\tfrac12-\eta)(t+h^2)\\
 &-\frac{a^4}{2}\{A^2t+3t^2+4Aht+6h^2t+h^4\},
 \qquad h=|u|.\notag
\end{align}
This cancellation is essential: an additive error of order $|u|/n$
would not be uniformly small compared with arbitrarily small $\lambda$.

\subsubsection{Finite normalization valid for every later index}
\label{dc:tail-normalization}
For clarity we give a finite recipe, rather than invoke an asymptotic
expansion. Set $N=1000$, $U=2/25$, $A=11/20$.
For a rational $z\ge0$ define
\[
 \sqrt z_-:=10^{-12}\lfloor10^{12}\sqrt z\rfloor,
 \qquad \sqrt z_+:=\sqrt z_-+10^{-12}.
\]
Integer square roots compute these directed bounds exactly. On a row
$[l,h]$, put $v_+=h(1-h)$ and use the fixed choices
\begin{gather*}
 \sigma_-=3.162277660,\quad\sigma_+=3.162277661,\qquad
 e=\max\{0,(4-\sigma_+)/2-h,\ l-(4-\sigma_-)/2\},\\
 P_0=(6+2\sigma_-)e^2/\sqrt{v_+}_+,\qquad
 \iota=6(1-\Theta)\sqrt{(1-\Theta)/v_+}_-(1-d_0),\\
 \beta=\min\{3(\sigma_-+1)B_-,3(\sigma_-+3)A_-\}-3P_0/2,\\
 \mathcal B=\min\{47D_{\rm row},\mathcal T/\kappa_0\},\quad
 \kappa_0=997/15000.
\end{gather*}
Here $\mathcal T$ is the original row's scalar contact resource. The
inequalities $\sigma_-^2<10<\sigma_+^2$ prove the distance bound used
in $P_0$; the central row containing $p_*$ gets $P_0=0$.
Apply \eqref{dc:noise-cap} with
$E=\mathcal B-NP_0$ and $n$ replaced by its lower bound $N$ in the
denominators. The original 45 central rows exclude indices
$j=0,\ldots,8$ and $j=28,\ldots,44$ by $E<0$.
For the 19 remaining rows, $p\in[.398,.436]$, this proves
$\lambda\le U$ and $d\le d_*:=\min(d_0,E/(N\iota))$ for every $n\ge N$.

On a retained row $[l,h]$, let $H_0=\max(H_-,H_+)$ and use
the following exact rational formulas:
\begin{gather*}
 \Delta=\max\{1,H_0\sqrt{h(1-h)+U/N}_++1/2\},\\
 r_-=l-\Delta/N,\quad r_+=h+\Delta/N,\quad v_0=r_-(1-r_-),\\
 Z=\frac{(1-2r_-)\Delta}{2v_0}+\frac1{12v_0}
                  +\frac{(1-2r_-)\Delta^3}{6Nv_0^2},\\
 B=\frac{U}{2v_-}+\frac{\Delta^2U}{2Nv_-^2}
                         +\frac{\Delta+1/4}{2v_-}+Z,
 \qquad v_-=l(1-l),\quad v_+=h(1-h).
\end{gather*}
The checked conditions are $0<r_-<r_+<1/2$, $Nr_->1$, $B<N$,
$h+\Delta<2$, and positivity of each denominator below. Finite Stirling
and Taylor bounds give
\begin{equation}\label{dc:stirling}
 \sqrt n\pi_k=\frac{\phi(0)}{\sqrt v}
    \exp\{-\delta^2/(2nv)+\zeta\},\qquad |\zeta|\le Z/n.
\end{equation}
For completeness, if $a_j=\log(j!)-(j+1/2)\log j+j$, then
\[
 0<a_j-a_{j+1}=\sum_{i\ge1}\frac{(2j+1)^{-2i}}{2i+1}
                       <\frac1{12j(j+1)}.
\]
Summing and using Wallis's product to identify
$\lim a_j=\log(2\pi)/2$ proves the remainder $0<r_j<1/(12j)$.
Substitution into the binomial mass and Taylor's theorem for relative
entropy supply \eqref{dc:stirling}. The three terms in $Z$ bound the
logarithmic prefactor, factorial remainder and entropy third derivative,
respectively. The latter derivative is $(2r-1)/(r^2(1-r)^2)$.

The exact phase derivative satisfies
\[
 \log f'=-\tfrac12\log(1+s/v)+\delta^2/(2nv)
                    -(\delta+u)^2/(2n(v+s))-\zeta.
\]
Hence $|\log f'|\le B/n$, and
$a_\ell=1-B/N$, $a_h=(1-B/N)^{-1}$ are valid.
The same formulas and the binomial recursion give
\begin{align*}
 b_k&\ge m:=\frac{.39886}{\sqrt{v_+}_+}
          \left(1-\frac{\Delta^2/(2v_-)+Z}{N}\right)>0,\\
 R_0&:=1+\frac{h+\Delta-1}{Nv_-[1-(\Delta-1)/(Nl)]},\qquad
 (1-\Delta/(Nl))\lambda\le t\le(1+\Delta/(Nl))\lambda.
\end{align*}
The ratio's first factor has the displayed bound and later rising factors
are at most one because $h+\Delta<2$. The variance identity is
$t=\lambda+\delta(s_1-s_0)$, with $|s_1-s_0|\le s/l$.

To bound the repair, the weaker consequence $G\ge3p-d$ of
\eqref{dc:cube-exact} gives
\[
 J=6l-3l^2-\tfrac32-d_*-\frac{3(1-2v_-)U}{8v_-N}>0.
\]
Using $C\ge\CE>.4$, the positive bracket, and the normal-variance
derivative already computed gives
\begin{align}\label{dc:tail-repair}
 \mathcal R&:=\sqrt n\{\Phi(\delta/\sqrt{nv})-\Phi(\delta/\sqrt{nw})\}
                      +C(\rho_B-\rho)\le\kappa\lambda/n,\\
 \kappa&=\frac{.4\Delta}{2v_-\sqrt{v_-}_-}
                       -\frac{.4J}{(v_++U/N)\sqrt{v_++U/N}_+}.\notag
\end{align}
The negative term is bounded in its correct direction. No positive-part
assumption is made about an arbitrary finite-cell cubic bracket.

\subsubsection{Four quartics, with their coefficients displayed}
The exact row arithmetic admits the following common envelope:
\[
 c\ge\frac{19987}{20000},\quad R_0\le\frac{20039}{20000},\quad
 B\le\frac{931}{250},\quad m\ge\frac{4011}{5000},\quad
 \kappa\le\frac{52}{125},\quad\Delta\le\frac{1067}{1000}.
\]
Use these constants in Lemma~\ref{dc:explicit-majorant} and
\eqref{dc:tail-loss}, and put $\omega_\pm=1\pm1067/398000$.
The function $\mathcal L(t,h)$ is concave in $t$, so its minimum is bounded
by $t=\omega_\pm\lambda$. Each endpoint expression minus $\lambda/20$
is concave in $\lambda\in[0,U]$. At zero noise it is nonnegative because
$1/2-\eta-a^2/8\ge9/25$. At $\lambda=U$, it is a quartic in
$h\in[0,1/2]$.
After $h=z/2$, the following integers divided by $10^6$ are lower bounds
for its five degree-four Bernstein coefficients:
\begin{center}
\begin{tabular}{ccrrrrr}\toprule
Phase&Variance&$b_0$&$b_1$&$b_2$&$b_3$&$b_4$\\\midrule
$u\ge0$&$\omega_-$&13449&2641&2242&12252&1883\\
$u\ge0$&$\omega_+$&13492&2626&2116&11962&1377\\
$u\le0$&$\omega_-$&14250&3114&2845&13442&3186\\
$u\le0$&$\omega_+$&14296&3100&2717&13145&2665\\\bottomrule
\end{tabular}
\end{center}
For $F(z)=\sum_{i=0}^4c_i z^i$, these coefficients are
$b_j=\sum_{i\le j}c_i\binom ji/\binom4i$; all displayed lower bounds
are positive. Thus $f(u)-\E P(W_k-u)\ge\lambda/20$ throughout the
closed domain. The equivalent score identity
\[
 \rho_B(C-C_B)=b_k\{Q/\pi_k-f(u)\}+\mathcal R
\]
and \eqref{dc:tail-repair} now give
\begin{equation}\label{dc:tail-final}
 \rho_B(C-C_B)\le-\left(\frac{4011}{100000}-\frac{52}{125000}\right)\lambda
                 =-\frac{19847}{500000}\lambda.
\end{equation}
For positive noise this contradicts $C\ge\CE$ and the binary bound
$C_B\le\CE$. At zero noise the actual law is Bernoulli and its actual
threshold is handled by the strict finite binary theorem, without dividing
by noise. This proves the central comparison for every $n\ge1000$ under
the original acquired central hypotheses.

\subsection{Numerical domains and the order of assembly}
\label{dc:assembly}
The analytic results above apply only after the same-law acquisition has
supplied their inputs. The numerical inequalities are summarized below.
The counts describe complete source domains, not samples.
\begin{center}
\begin{tabular}{>{\raggedright\arraybackslash}p{.22\linewidth}>{\raggedright\arraybackslash}p{.28\linewidth}>{\raggedright\arraybackslash}p{.32\linewidth}}\toprule
Branch&Domain&Certificate test\\\midrule
Uniform left&39 raw rows, 3385 retained actual-index pairs;
 $47\le n\le326$&Global quartic and 24 moment tests per cell.\\
Uniform central, finite&45 raw rows, 25908 retained pairs;
 $47\le n\le999$&Finite binary reserve and 24 moment tests per cell.\\
Uniform central, tail&45 raw rows, every $n\ge1000$&Scalar bounds and four displayed quartics.\\
Earlier finite families&737 row/index pairs in the table below&
 Global quartic and 48 signed moment tests per cell.\\\bottomrule
\end{tabular}
\end{center}

The uniform-left upper indices follow from a positive
Bernoulli-defect lower bound on each row. The original conservative bound
uses $(1-\Theta)^{3/2}\inf D_B$ in $nD\le47D_{\rm row}$; the largest
retained integer is $326$. The strengthened common budget
\eqref{dc:linear-budget} then sharpens noise within each retained cell.
The central finite census uses
$nP_0\le\min(47D_{\rm row},T_H/\kappa_0)$ with
$\kappa_0=997/15000$ and the unrefined scalar resource when computing its
global integer cap. Row $19$ contains $p_*$ and has $P_0=0$; it retains
every integer $47$ through $999$. Indices above a positive cap are excluded by
the defect inequality, not by an absent proof object.

\begin{proposition}[Earlier finite comparison]
\label{j:local:direct-finite-numerical}
\label{dc:earlier-numerical}
On the original acquired rows in the following table, complete closed
count/probability/phase covers satisfy Theorem~\ref{dc:finite-test}, using
the stated actual index and target:
\begin{center}
\begin{tabular}{lcccc}\toprule
Family&Rows&Actual $n$&Target&Outward strip\\\midrule
Positive $39$&78&39&$\CE+741/(4000\cdot39)$&$.55$ or $.575$\\
Central $40$&48&40,41,42&$\CE$&$.55$\\
Left $40$&49&40,41,42&$\CE$&$.55,.575$ or $.6$\\
Central $42$&47&43,44,45,46&$\CE$&$.55$\\
Left $42$&45&43,44,45,46&$\CE$&$.575$\\\bottomrule
\end{tabular}
\end{center}
The complete exact acceptance checks all 737 pairs, using 2090 quartics
and 1078 resource exclusions. Its 48-corner tests retain both signs of
the cubic correction; no subdivision or unresolved cell remains.
\end{proposition}

Contact proves each coordinate strip in this table before the comparison. The floor $40$
acquisition applies at actual index $40$ with predecessor decrement at most .19,
or at actual indices $41,42$ with decrement zero; the indexed central extension is
retained as a separate proof. The floor $42$ acquisition preserves
$nD\le42D_{\rm row}$ and its original finite-band support theorem.
At actual indices $43$ through $46$ the variable $n$ is retained in every mass, variance,
score reserve and normal correction.

The positive-$39$ acquisition has $K=741/4000$, predecessor allowance .33
and its original positive charge. Its target is $\CE+K/39$, both in
$\rho_B(T-C_B)$ and in the two cubic multipliers. Its already acquired
defect bound retains the $K$ charge; omitting an additional negative
$K\rho$ term from a subsequent phase-resource refinement merely weakens
that refinement. It must not be mistaken for a new zero-allowance
acquisition. This prerequisite is proved before the sharp-$39$ input is
used; sharp-$39$ is never assumed to obtain its own entry bound.

\begin{proposition}[Central finite comparison]
\label{dc:central-numerical}
Every one of the $25\,908$ required central row/index pairs has a freshly
verified finite binary reserve and a complete exact cover satisfying
Theorem~\ref{dc:finite-test}. Missing, duplicate or extraneous pairs
reject. The full binary-gap checks and the separate exact comparison
replay cover the complete declared domain.
\end{proposition}

\paragraph{Assembly of the comparison.}
\label{j:local:comparison}
The earlier finite induction uses Proposition~\ref{dc:earlier-numerical} for its positive-$39$ prerequisite
and actual indices $40$--$42$ comparisons. The original induction through $39$ and its
analytic moment/support licenses otherwise remain unchanged.
The actual indices $43$--$46$ local exclusion uses the floor $42$ rows above.
The uniform entry assumptions place a hypothetical maximizer into an
original left or central row. The left comparison, the completed central
finite assertion, and \eqref{dc:tail-final} then exclude all $n\ge47$.
The original initialization and complements cover the remaining
probabilities; the comparison does not silently replace those covers.

Equality at an
index $n\ge47$ is handled only after the weak all-index bound has been
proved: both deleted bounds are then at most $\CE$, the excess and
decrement are zero, and ratio attainment supplies a maximizing law.
The uniform acquisition includes that boundary. Positive noise is excluded
by the strict comparisons, and zero noise by the strict finite binary
theorem. This explains why pointwise strictness alone would not be enough
to prove strictness of the supremum.

\begin{proposition}[Uniform exclusion]
\label{j:local:uniform}
No attained positive closed-event ratio maximum
satisfies
\[
 n\ge47,\quad 0\le n(C_n-\CE)\le.7,\quad
 n(C_{n-1}-C_n)_+\le.19,\quad C_n<.4247,
 \quad C_{n-1},C_{n-2}<.4253.
\]
\end{proposition}
\begin{proof}
The retained uniform acquisition places the same law in a raw left or
central row. The three uniform branches in the status table cover every
index and every phase/noise value. Positive noise is excluded by their
strict score inequalities. At zero noise use the binary theorem at the
actual threshold. The original acquisition uses the excess only as a
nonnegative resource term, so this argument includes excess zero.
\end{proof}

\begin{proposition}[Exclusion at indices forty-three through forty-six]
\label{j:local:finite4346}
No attained positive ratio maximum at an integer $43\le n\le46$ satisfies
\[
 0<n(C_n-\CE)\le.54,\quad n(C_{n-1}-C_n)_+\le.19,\quad
 C_n<.4226,\qquad C_{n-1},C_{n-2}<.4253.
\]
\end{proposition}
\begin{proof}
The original floor-$42$ initialization, support and acquisition cover
every probability panel under these hypotheses. Their surviving central
and left rows retain $nD\le42D_{\rm row}$ and their independently proved
coordinate strips. Apply the corresponding actual-index part of
Proposition~\ref{dc:earlier-numerical}, with target $T=\CE$.
Its count, phase and noise cover contains the same maximizing law and
gives a strict score below $\CE$, a contradiction.
\end{proof}

\begin{corollary}[Equality exclusion]
\label{j:local:equality}
Once $C_j\le\CE$ is proved at all indices, one has $C_n<\CE$ for every $n\ge47$.
\end{corollary}
\begin{proof}
If equality held, attainment and the already proved weak bounds would
give all hypotheses of Proposition~\ref{j:local:uniform}, with zero
excess and zero predecessor decrement. This contradicts that proposition.
\end{proof}
